\documentclass[11pt,letterpaper,reqno]{amsart}
% BEGIN INLINE manuscript-style
\makeatletter
\NeedsTeXFormat{LaTeX2e}[2020/10/01]
% Use with \documentclass[11pt,letterpaper,reqno]{amsart}.
% This file controls presentation only. It supplies no mathematical content.

\RequirePackage[T1]{fontenc}
\RequirePackage{lmodern}
\RequirePackage{microtype}
\RequirePackage[letterpaper,margin=1in,headheight=15pt]{geometry}
\RequirePackage{amsmath,amsthm,amssymb,mathtools}
\RequirePackage{mathrsfs}
\RequirePackage{enumitem}
\RequirePackage{perpage}
\RequirePackage{longtable}
\RequirePackage{xcolor}
\definecolor{manuscriptred}{rgb}{0.6,0.1,0.1}
\definecolor{manuscriptblue}{rgb}{0.1,0.3,0.5}
\definecolor{manuscriptmagenta}{HTML}{60445C}
\RequirePackage[breaklinks=true,colorlinks=true,
  linkcolor=manuscriptred,citecolor=manuscriptblue,urlcolor=manuscriptmagenta]{hyperref}
\RequirePackage{cleveref}
\hypersetup{pdfauthor={}}

% Retain the AMS paragraph, title, heading and running-head conventions.
% Equation counters restart each section: 2.i, 2.ii, ... .
% Appendix sections keep their letters: A.i, A.ii, ... .
\numberwithin{equation}{section}
\renewcommand{\theequation}{\thesection.\roman{equation}}
% Prose references include "eq." and omit parentheses around the number.
\renewcommand{\eqref}[1]{\textup{eq.~\ref{#1}}}
\crefname{equation}{eq.}{eqs.}
\Crefname{equation}{eq.}{eqs.}
\creflabelformat{equation}{#2#1#3}
\setcounter{tocdepth}{1}
\renewcommand{\thefigure}{\Roman{figure}}
\renewcommand{\thetable}{\Roman{table}}
\newcommand{\sineedspace}[1]{\par\begingroup
  \dimen0=#1\relax
  \vskip0pt plus\dimen0\penalty-100
  \vskip0pt plus-\dimen0\vskip\dimen0
  \penalty9999\vskip-\dimen0\vskip0pt\endgroup}
\MakePerPage{footnote}
\renewcommand{\thefootnote}{%
  \ifnum\value{footnote}>9\arabic{footnote}\else\fnsymbol{footnote}\fi}

% Upright statements. Main results are A, B, C; supporting statements
% share a section-local sequence. These end marks indicate typography,
% never proof completeness, validation, or a deferred-proof status.
\theoremstyle{definition}
\newtheorem{ms@theorem}{Theorem}
\renewcommand{\thems@theorem}{\Alph{ms@theorem}}
\newtheorem*{ms@theoremstar}{Theorem}
\newtheorem*{ms@theoremzero}{Theorem 0}
\newtheorem{ms@proposition}{Proposition}[section]
\newtheorem{ms@lemma}[ms@proposition]{Lemma}
\newtheorem{ms@definition}[ms@proposition]{Definition}
\newtheorem{ms@corollary}[ms@proposition]{Corollary}
\newtheorem*{ms@corollarystar}{Corollary}

\newcommand{\ms@majorstart}{%
  \pushQED{\qed}\renewcommand{\qedsymbol}{$\blacksquare\square$}}
\newcommand{\ms@minorstart}{%
  \pushQED{\qed}\renewcommand{\qedsymbol}{$\blacksquare$}}
\newenvironment{theorem}[1][]{\ms@majorstart\ms@theorem[#1]}
  {\popQED\endms@theorem}
\newenvironment{theorem*}[1][]{\ms@majorstart\ms@theoremstar[#1]}
  {\popQED\endms@theoremstar}
\newenvironment{theoremzero}[1][]{\ms@majorstart\ms@theoremzero[#1]}
  {\popQED\endms@theoremzero}
\newenvironment{proposition}[1][]{\ms@minorstart\ms@proposition[#1]}
  {\popQED\endms@proposition}
\newenvironment{lemma}[1][]{\ms@minorstart\ms@lemma[#1]}
  {\popQED\endms@lemma}
\newenvironment{definition}[1][]{\ms@minorstart\ms@definition[#1]}
  {\popQED\endms@definition}
\newenvironment{corollary}[1][]{\ms@majorstart\ms@corollary[#1]}
  {\popQED\endms@corollary}
\newenvironment{corollary*}[1][]{\ms@majorstart\ms@corollarystar[#1]}
  {\popQED\endms@corollarystar}
% The source's lemma/proposition variants are deliberately not carried over:
% choose one marker for each statement type rather than encode proof status.

\theoremstyle{remark}
\newtheorem{remark}[ms@proposition]{Remark}
\newtheorem*{remark*}{Remark}
% The ordinary amsthm proof environment retains its open square.

\crefname{ms@theorem}{theorem}{theorems}
\Crefname{ms@theorem}{Theorem}{Theorems}
\crefname{ms@proposition}{proposition}{propositions}
\Crefname{ms@proposition}{Proposition}{Propositions}
\crefname{ms@lemma}{lemma}{lemmas}
\Crefname{ms@lemma}{Lemma}{Lemmas}
\crefname{ms@definition}{definition}{definitions}
\Crefname{ms@definition}{Definition}{Definitions}
\crefname{ms@corollary}{corollary}{corollaries}
\Crefname{ms@corollary}{Corollary}{Corollaries}
\crefname{remark}{remark}{remarks}
\Crefname{remark}{Remark}{Remarks}

% Bibliography configuration for this manuscript.
\makeatother
% END INLINE manuscript-style

\title[Scalar sources and uniform insertions]{Scalar sources and uniform insertions for lattice Yang--Mills fields}
\date{Working draft, 9 November 2026}
\begin{document}
\raggedbottom
\begin{abstract}
We develop scalar action-density sources on admitted lattice gauge
trajectories, with the real reference construction held fixed. Local
source formation and retained-integral estimates give holomorphic
source energies and a complex partition upper bound. Reflection and
chessboard estimates transfer these bounds to every periodic Wilson
Gibbs state at fixed cutoff, using a common cell normalizer independent
of the state. Rooted domain counting yields graded separated moments.
A separate argument on the reflected observable algebra gives a
one-field insertion bound independent of the surrounding polynomial
degrees and number of fields. The local analytic formation and
history bounds remain explicit inputs. Conditional passage of the
insertion estimate to convergent separated moments is included;
existence of a continuum limit is not asserted.
\end{abstract}
\maketitle
\tableofcontents

\section{Introduction}

A spatially varying source for the Wilson action must pass through
the same local transformations as the original measure. At complex
source values positivity is unavailable. The source dependence of
the local interactions, their extracted coefficients and their
numerical vacuum terms must therefore be controlled before summing
the histories. An unsourced partition bound does not provide that
control. We keep the real reference data fixed, identify the exact
source transport and local scalar subtractions, and deduce the
sourced upper bound from the stated local analytic and integral
estimates.

The passage to infinite volume uses a different argument. Reflection
and chessboard domination compare every periodic Gibbs state with
finite tori at fixed lattice spacing. The holomorphic normalizer
used in this comparison is selected from the finite-torus local
energies before a state is chosen. Its first derivative may grow
with the cutoff, but it cancels when comparing state means or
centering observables at another state. This permits estimates for
all such states without treating each state as a torus limit.

There are two distinct conclusions for products of fields. A rooted
domain count and a choice of cell offsets give a product source
polydisc for separated pieces. Cauchy's formula then yields one
fixed factorial order for moments and cumulants. To control an
insertion in the reflected norm of an arbitrary polynomial, we
instead repeat a cell insertion along a time column. Log-convexity
of its raw reflection moments removes the endpoint constants.
Boundedness on the Hilbert quotient is a consequence of this
argument, rather than a premise of it.

The localization and connected expansion methods have their
antecedents in Ba\l aban's small-field construction
\cite[\S\S1--2, pp.~2--22]{Bal88c}. The present deductions also use
the finite reflection and Gibbs-state results of \cite{LP}, the
local block conventions of \cite{BA}, the regular-family interfaces
of \cite{EM}, and the admitted real trajectories and unsourced floor
of \cite{RT}. The actual nonlinear source domains, local formation
and retained mass estimates are named hypotheses at the points
where they enter. The conclusions do not include an interacting
two-point asymptotic law or a construction of continuum fields.

\section{Finite sources and fixed reference data}
\label{si:sec-setting}

Fix an odd integer $L\ge3$ and a cutoff $K\ge0$. Put $a=L^{-K}$.
Fix also a nonnegative integer $m_0$ and put $M=L^{m_0}$;
localization cubes have side $M$ in their own mesh units.
On the periodic lattice $a\mathbb Z^4/\ell\mathbb Z^4$, assign
$U_b\in\operatorname{SU}(2)$ to every positive bond and set
$U_{\bar b}=U_b^{-1}$. All periods and localization partitions used
below are required to be compatible. Each unoriented plaquette is
counted once, with a prescribed orientation and barycentre $x(p)$.
Write
\begin{equation}\label{si:eq-finite-source}
 a_p(U)=1-\tfrac12\operatorname{Tr}U_p,\qquad
 A(U)=\sum_p a_p(U),\qquad
 \mathcal O_a(f)=\sum_p f(x(p))a_p(U).
\end{equation}
Here $0\le a_p\le2$ on real configurations. There is no additional
factor $a^4$ in the definition of $\mathcal O_a$.

The half trace fixes the smooth-field normalization. For a
Hermitian traceless matrix $F$ use $|F|^2=\tfrac12\operatorname{Tr}F^2$.
If on a fixed compact region the elementary holonomies of a
smooth connection have logarithms
$i a^2 F_{\mu\nu}(x(p))+O(a^3)$ with uniform remainder, these
small logarithms are traceless. The quadratic term gives
\[
 a_p(U)=\tfrac12 a^4 |F_{\mu\nu}(x(p))|^2+O(a^5).
\]
For a fixed continuous test $f$ supported in this region, the weighted sum in
\eqref{si:eq-finite-source} is a Riemann sum for
$\tfrac12\int f\sum_{\mu<\nu}|F_{\mu\nu}|^2$. The estimates below
apply to the lattice Wilson measure through the finite observable
in \eqref{si:eq-finite-source}.

Under a site gauge the plaquette product is conjugated at its
base point, so its trace and every $\mathcal O_a(f)$ are gauge
invariant. Reversing a plaquette orientation replaces its matrix
by its inverse, which has the same trace in $\operatorname{SU}(2)$.
A lattice translation or signed coordinate permutation carries
the observable to the one with the correspondingly transported
barycentric test. Reindexing the bonds preserves product Haar
measure and the Wilson action, so finite-volume expectations have
this covariance whenever the transformation preserves the torus.
With Haar probability measure on each bond, set
\[
 Z_{a,\ell,\beta}(W)
   =\int \exp\{-\beta A(U)+\mathcal O_a(W)\}\prod_b dU_b,
 \qquad \beta>0.
\]
The unsourced probability law is denoted by $\mu_{a,\ell,\beta}$.
Its denominator $Z_{a,\ell,\beta}(0)$ is strictly positive.

The source data are a finite triple $w=(c,(f_i),(z_i))$, whose
function is $W=c+\sum_i z_i f_i$. Here $c,z_i\in\mathbb C$ and
the $f_i$ are real Lipschitz torus functions; the family may be
empty. Fix the geometric constant $c_M>0$. For $s\ge0$, put
\begin{equation}\label{si:eq-source-norm}
 \|f\|_{\sharp,s}=\max\{\|f\|_\infty,c_M L^{-s}\operatorname{Lip}(f)\},
 \qquad
 \|w\|_{\flat,s}=|c|+\sum_i|z_i|\|f_i\|_{\sharp,s}.
\end{equation}
At $s=0$ omit the second subscript. This is the physical source
size used in the volume-limit estimates.
For real $C^1$ functions the Lipschitz seminorm can be replaced by
the supremum of the Euclidean gradient norm. For a complex
Lipschitz test use the same $\sharp,s$ formula, with the
Lipschitz seminorm of its map into $\mathbb C$. Its real and
imaginary parts each have norm at most that of the test, so
their two-profile representation costs at most twice that norm.
The number
$\|w\|_{\flat,s}$ belongs to the chosen triple; it can change when
the same function is represented by another list. It bounds both
$\|W\|_\infty$ and $c_M L^{-s}\operatorname{Lip}(W)$.
Every local source term must agree for two representations with
the same values on all the source points that it reads.
Holomorphy means holomorphy in $(c,z)$ for each fixed real family,
with these compatible values. The ordinary source class consists
of all triples with $\|w\|_\flat<2r$, for one fixed $r>0$.

Write $a_s=L^{-s}$ for physical depth scales. If $m=K-s$ and
$j\le m$, the source size gives the exact freezing estimate
\begin{equation}\label{si:eq-source-freezing-scale}
 c_M a_{K-j}\operatorname{Lip}(W)
   =L^{-(m-j)}c_M a_s\operatorname{Lip}(W)
   \le L^{-(m-j)}\|w\|_{\flat,s}.
\end{equation}
This gain concerns variation of the same source function across
an earlier physical block.

The real functions permitted in this ball include global torus
profiles obtained by reflection of a function supported inside one
cell. They need not have support in a single ball. This inclusion
will allow the chessboard construction to use one real profile for
all even copies and another for all odd copies.

\begin{samepage}
\begin{lemma}[Finite source derivatives]\label{si:lem-finite-derivatives}
For a fixed finite torus and real bounded tests $f_1,\ldots,f_n$, the
function $Z(\sum_i z_if_i)$ is entire in $z\in\mathbb C^n$. Its
mixed derivatives are obtained by inserting the corresponding
products of $\mathcal O_a(f_i)$ in the integral. A holomorphic
logarithm of $Z(z)/Z(0)$ exists near zero. Its mixed derivative,
with one derivative in each variable, is
\begin{equation}\label{si:eq-cumulant-partitions}
 \kappa_n^{\mu}(\mathcal O_a(f_1),\ldots,\mathcal O_a(f_n))
 =\sum_{\pi}(-1)^{|\pi|-1}(|\pi|-1)!
     \prod_{B\in\pi}\mu\!\left(\prod_{i\in B}\mathcal O_a(f_i)\right),
\end{equation}
where $\pi$ ranges over the set partitions of $\{1,\ldots,n\}$.
For $n\ge2$ this value is unchanged if any test observable is
shifted by a constant. Conversely, each product moment is the
sum over set partitions of the products of the corresponding
block cumulants.
\end{lemma}
\end{samepage}
\begin{proof}
Each observable is bounded at this fixed cutoff and volume. On a
compact parameter set, its exponential and every fixed derivative
have one integrable constant bound times the positive Wilson
density. Dominated convergence proves differentiation under the
integral and holomorphy. Continuity and $Z(0)>0$ give a neighborhood
on which $Z(z)/Z(0)$ lies in the disc of radius $1/2$ about one.
The logarithm series is holomorphic there.

Put $F(z)=Z(z)/Z(0)$. In
$\log F=\sum_{k\ge1}(-1)^{k-1}(F-1)^k/k$, the coefficient using
each of the $n$ variables once receives contributions only from
$k\le n$. Assigning variables to the $k$ nonempty ordered factors
amounts to ordering the blocks of a set partition. The $k!$
orderings and the factor $1/k$ give $(k-1)!$, proving
\eqref{si:eq-cumulant-partitions}. Expanding $F=\exp(\log F)$
gives the inverse formula: each ordering of the $k$ blocks is
cancelled by the exponential coefficient $1/k!$. Shifting an observable multiplies
$F$ by the exponential of a linear function, which changes only
the first derivatives of its logarithm.
\end{proof}

We use an admitted real reference run from the trajectory
construction \cite{RT}. Its chosen weights $x_j$, reference
weights $\theta_j$, cutoffs, partitions, gauges and positive history
kernels are frozen when a source is introduced. The reference
Gaussian is the one of that real run. A complex source changes the
inserted density; it supplies no positivity for a new complex
Wilson law. The local analytic formation and retained-integral
bounds assumed below concern this fixed real construction.

For estimates on infinite lattices, a periodic Wilson Gibbs state
means a probability measure invariant under a finite-index group
of lattice translations and satisfying the Wilson DLR
specification at the same fixed $a$ and $\beta$. We use the
reflection, chessboard and all-state domination results of
\cite{LP}. In particular \cite[Theorem~B]{LP} gives invariance
under every lattice translation for these states. The volume
limit is taken at fixed $a$ and $\beta$.
Statements comparing two states use the same cutoff and coupling
for both. Existence of a continuum limit is not an input to these
finite-spacing arguments.

Constants denoted by $C$ may increase between estimates. Unless a
dependence is displayed or stated, they depend only on the fixed
geometric and analytic input bounds, not on the cutoff, compatible
volume, source representation or Gibbs state. Constants indexed by
the number of observables may depend on that number. Fixed-cutoff
bounds for point counterterms and for normalizers before taking a
volume subsequence are explicitly distinguished from uniform bounds.


\section{Source transport with fixed reference data}\label{si:sec-frozen-transport}

At level $j$ the mesh is $h_j=L^{j-K}$. A source is evaluated at the
bare plaquette barycentres, including the plaquette points used as
anchors by earlier levels. Fix the reference choices before assigning
a source. These choices include the background and current maps, real
Gaussian precision and gauge weight, every test and guard, and the
reference laws used when an interior field is replaced. The sourced
Wilson weight is $e^{-A(\beta-W,U)}$, where
$A(u,U)=\sum_p u(p)a_p(U)$. In particular a constant real source $v$
changes the bare weight to $\beta-v$ in this same reference.

\begin{samepage}
\begin{proposition}[Transport of the bare source]\label{si:prop-frozen-transport}
Use a finite admitted prefix of the fixed reference run. Let
$K_h(U,dV)$ be its positive history kernels through any level $j$,
with $\sum_h K_h(U,\mathcal V_j)=1$. Then the sourced terms are
\begin{equation}\label{si:eq-frozen-transport}
 \tau_h^W(dV)=\int e^{-\beta A(U)}e^{A(W,U)}K_h(U,dV)\,dU,
 \qquad \sum_h\tau_h^W(\mathcal V_j)=Z(W).
\end{equation}
On each finite source family the measures $\tau_h^W$ are entire in
total variation. For a real source they are positive measures. The
same formula transports a bounded Borel gauge-invariant insertion
$F(U)$ by putting $F$ in the integrand.
\end{proposition}
\end{samepage}
\begin{proof}
The history kernels are formed by the normalized guarded averages,
gauge draws, coordinate changes and reference replacements of the
finite run \cite{RT}. Each actual move acts linearly on the incoming
measure, and the outgoing kernels sum to a probability kernel. A
gauge draw uses its normalized kernel, with the orbit representatives
and every test retained in their order. Integrating the draw against
its input gives the gauge-weight formula when that input is invariant
under the gauge subgroup of the draw. For this identification use the actual retained-gauge equivariance
at each draw, which preserves invariance of the weighted incoming
term under that draw subgroup. The bare Wilson source and the
allowed insertion are gauge invariant. Composition
and integration give \eqref{si:eq-frozen-transport}. On bounded
sets of source coefficients the exponential series converges uniformly
on the compact bare configuration space, since $0\le a_p\le2$ and
there are finitely many plaquettes. Integration against any fixed
finite history kernel is bounded in total variation. This proves the
entire assertion and the same argument gives the insertion formula.
For real profiles the integrand defining the sourced measure is
nonnegative. The partition of outcomes depends on the fields and
the fixed tests, so no outcome is introduced by complexification.
\end{proof}

\begin{samepage}
\begin{proposition}[Block observables outside replaced interiors]
\label{si:prop-block-observables}
Suppose the actual averaged stage sends each incoming history term
under $\mathcal M$ to the sum of its tested continuations. Suppose also
that for each tested continuation the replacement stage preserves the
marginal on bonds exterior to the designated replacement components.
It suffices that these identities hold on gauge-invariant bounded Borel
tests, including all such tests of the exterior bonds in the second
identity. Use the covariant local block map of \cite{BA}. At $W=0$
there is a joint probability law of $U,h,V$ whose bare marginal is the
Wilson probability and whose history/block marginal is
$Z(0)^{-1}\sum_h\delta_h\otimes\overline{\tau_h^0}$, where the bar is
full site-gauge averaging. The law is invariant under simultaneous
bare gauges and their restrictions on $V$. On a level-$j$ bond whose
complete backward block read misses the enlarged components acted on
along $h$, it has $V(c)=\mathcal M^jU(c)$ almost surely.

For bounded gauge-invariant $F$, let $F_h(V)$ be the invariant version
of the conditional expectation of $F(\mathcal M^j U)$ given $(h,V)$.
Then
\begin{equation}\label{si:eq-block-observable}
 \mu(F\circ\mathcal M^j)=Z(0)^{-1}\sum_h\int F_h\,d\tau_h^0,
 \qquad |F_h|\le\|F\|_\infty\quad\tau_h^0\text{-almost everywhere}.
\end{equation}
If $F$ reads only the bonds just specified for $h$, then $F_h=F$
$\tau_h^0$-almost everywhere. Omitting any event of histories in
\eqref{si:eq-block-observable} costs at most its probability times
$\|F\|_\infty$.
\end{proposition}
\end{samepage}
\begin{proof}
Average the incoming, tested and completed terms under the full site
gauge groups. The identities on invariant tests become measure
identities, since any bounded test of an invariant measure has the
same integral as its Haar average. Projection to the exterior
commutes with gauge averaging. The block map commutes with the
restriction of a site gauge, and restriction carries Haar probability
to Haar probability on the coarser sites.

Fix an incoming history $h$. If $\lambda$ is its averaged term and
$\lambda_b$ the averaged tested terms for its continuations, then
$\sum_b\lambda_b=\mathcal M_*\lambda$. Write
$r_b=d\lambda_b/d(\mathcal M_*\lambda)$. From $V$ with measure
$\lambda$ choose $b$ with probability $r_b(\mathcal M V)$ and retain
$\mathcal M V$ as the tested field. For each $b$ the sum of completed
terms, with their continuation labels, has the same exterior
marginal as $\lambda_b$. Disintegrate this sum over the exterior
field and use its conditional probability as the second draw. The
field after this draw agrees with the tested field on the exterior.
All configuration spaces are finite products of compact groups and
all label sets are finite, so the conditional laws exist. Iterating
these two draws gives the required marginals. At each step the complete block read includes the fine plaquettes
for its branch tests and all chosen contours, including projected
short seam reads. It lies in the projected double block by the
locality result of \cite{BA}. Under nested centred blocks all its
recursive reads lie in the backward double block. Consequently, if the complete backward read misses all the
acted-on enlarged components, induction identifies every read with
the corresponding read of the bare iterated map. The final read
therefore agrees. Averaging the joint law under the simultaneous
gauge action preserves both marginals and these equality events.

The conditional expectation can be averaged under site gauges on
$V$. Invariance of the joint law and surjectivity of gauge
restriction make the average a version of the same conditional
expectation. Its essential bounds and its equality to $F$ on the
specified histories concern invariant events and so pass from
$\overline{\tau_h^0}$ to $\tau_h^0$. Integrating the conditional
expectation proves \eqref{si:eq-block-observable}. Positivity of the
omitted measures proves the last bound. In particular a supplied
one-cell rarity bound gives the corresponding finite union bound
for a list of cells.
\end{proof}

\section{The running scalar weight}\label{si:sec-running-profiles}

On each fixed formation window, let $h_y^{(j)}(p)$ be the nonnegative
tents of the step with $\sum_y h_y^{(j)}(p)=1$. The anchor $p_y$ is a
fixed bare plaquette and $h_y^{(j)}(p)\ne0$ implies
$d(x(p),x(p_y))\le C M h_j$. Let $0\le\varphi_j(p)\le1$ be the
actual stage profile. Write $G_{j,y}(v)$ for the constant-source
whole-lattice member used at $y$, with the fixed read and pointing
of the step. It includes the companion needed for its physical
quadratic readout. For constant $v$ the member has the covariance
required by the pointed coefficient theorem of \cite{EM}. At the
trivial pair $P_\circ=(\mathbf1,0)$ subtract its vacuum value, and
use its own coefficient $b_j(v)$ to subtract its local action.
In the unit-Wilson convention this is the operation
\begin{equation}\label{si:eq-member-subtraction}
 G_{j,y}(v)-G_{j,y}(v)(P_\circ)
                 -b_j(v)A(h_y^{(j)},\cdot).
\end{equation}
Here the action denotes the actual pointed classical member, with
its exact restriction and inherited transport. At a site where the
array selects $G_{j,y}(W(p_y))$, the coefficient in this formula is
$b_j(W(p_y))$. The coefficient cancels the quadratic moment of that
member by the linear readout in \cite{EM}. There is no requirement
that the array of site-dependent members be translation invariant.
The Euclidean members are supplied on the common reference and
common formation domains. When read on an admitted torus they use
those same coefficients. A different torus member with no common
restriction would require its own readout.

The complete old pointed member and its tent action use a single
window read and a single bracket functional at fixed $j,y$. Linearity
allows the bracket to be applied to their vacuum- and
quadratic-subtracted sum before estimating it. For an infinite old
domain sum, normal convergence on the switch disc justifies this
identity by Cauchy's formula. Truncation to near domains leaves the
omitted part in the separate far-tail estimate. The coefficient is
attached to the complete old pointed member. It is neither an
individual domain coefficient nor a claim about the readout of
the new bracket output. Differentiating the fixed linear readout
and subtraction in \eqref{si:eq-member-subtraction} gives the same
subtraction for the source derivative with its own coefficient.

\begin{samepage}
\begin{proposition}[Scalar shifts and their source range]
\label{si:prop-source-profiles}
Suppose the actual constant-source formation is available at each
step on the admitted half product of the relative weight discs of
\cite{RT}, with common older reads. At step $i$ let the nonconstant
unmarked coefficient be holomorphic on the full relative product,
with bound $C\theta_{i-1}^{-1/2}(\log\theta_{i-1})^{\mathfrak p_0}$,
and let the formed marked coefficient have bound
$C\theta_{i-1}^{-p_R}$ on $|v|<4r$, where $p_R>1$. The marked
hypothesis is for the actual holomorphic sourced term including inherited
remainders. Both coefficient maps are real on real data.
Assume the real deviations at the used levels lie in
the quarter products and that the relative half products have
at least $6r$ of room beyond them. Suppose
\[
 \theta_j\ge H,\qquad
 \theta_i-\theta_j\ge\lambda(j-i)-C_0\quad(i\le j),
 \qquad\lambda>0.
\]
For $H$ sufficiently large at fixed $r$, the successive coefficients
on $|v|<4r$ have radial bounds
\begin{equation}\label{si:eq-source-slopes}
 d_i=\sup_{0<|v|<4r}\frac{|b_i(v)-b_i(0)|}{|v|},\qquad
 \sum_{i\le j}d_i\le C\{\theta_j^{-1/2}(\log\theta_j)^{\mathfrak p_0}
                           +r^{-1}\theta_j^{1-p_R}\}\le\tfrac16.
\end{equation}
The bound is uniform in admitted volume and prefix length. For the
ordinary variable source define
\begin{equation}\label{si:eq-running-shift}
 \begin{split}
 c_i(p;W)&=\varphi_i(p)\sum_y h_y^{(i)}(p)b_i(W(p_y)),\\
 \zeta_j(p;W)&=W(p)+\sum_{i\le j}\{c_i(p;W)-b_i(0)\varphi_i(p)\},\\
 x_j(p)-\zeta_j(p;W)&=\beta-W(p)-\sum_{i\le j}c_i(p;W).
 \end{split}
\end{equation}
This is the exact Wilson weight after the local action subtractions.
Each value is holomorphic in the source coefficients and depends
only on bare source values within distance $CMh_j$ of $p$. In
particular $|\zeta_j(p;W)|\le(4/3)\|W\|_\infty$. At a constant
real source $|v|<r$ the uniform interior scalar shift satisfies
$2/3\le\zeta_j'(v)\le4/3$, and at zero its derivative differs from
one by at most the sum in \eqref{si:eq-source-slopes}.
\end{proposition}
\end{samepage}
\begin{proof}
Write $z_l=\theta_l-x_l$ in the constant-member reference run.
The fixed-reference coefficient is its common one-loop term
plus the ordinary analytic coefficient on
$(z_l+\zeta_l(v))_{l<i}$ and the actual marked coefficient
$q_i(v)$. This is the prescribed regular and marked splitting
of the formed coefficient. At every half product the absolute derivative row
bound of the relative-disc estimate in \cite{RT} is at most
$C\theta_{i-1}^{-3/2}(\log\theta_{i-1})^{\mathfrak p_0}$ at step $i$.
Indeed $\theta_l\ge\theta_{i-1}-C_0$ for $l<i$. If the preceding
radial sum is at most $1/6$, then $|\zeta_l(v)|\le(4/3)|v|<6r$,
so the entire segment for the next coefficient lies in the half
product. The formed marked term $q_i(v)$ satisfies
$|q_i(v)-q_i(0)|\le C r^{-1}\theta_{i-1}^{-p_R}|v|$ by Schwarz's
lemma on its actual source disc. The derivative row estimate and
integration on the unmarked segment thus bound $d_i$ by the sum of
these two envelopes with a fixed constant. Using
$\theta_{i-1}\ge\theta_j+\lambda(j-i+1)-C_0$ and comparing with
integrals gives \eqref{si:eq-source-slopes}. The sum over every
prefix is at most $1/6$ when $H$ is large enough. Induction starts at
$\zeta_0=v$ and admits each next weight before forming the next
coefficient. This proof concerns the analytic coefficient box, with
the separate real weight tube imposed on the physical formation.

In \eqref{si:eq-running-shift}, expand the real weight
$x_j(p)=\beta-\sum_{i\le j}b_i(0)\varphi_i(p)$ and use the partition
of unity. At each step the difference of successive shifted weights
is exactly $c_i(p;W)$. Linearity of the action therefore compensates
\eqref{si:eq-member-subtraction} at its own anchor. The distance
claim follows from the support of the tents and $h_i\le h_j$ for
$i\le j$. Positivity of the tents and \eqref{si:eq-source-slopes}
give the shift bound. For each holomorphic quotient
$(b_i(v)-b_i(0))/v$, the bound is $d_i$ on $|v|<4r$. At $|v|<2r$,
Cauchy's estimate gives $|b_i'(v)|\le2d_i$. In the constant interior
all stage factors are one, and reality gives the stated derivative
bounds. At zero use $|b_i'(0)|\le d_i$.
\end{proof}

For an interior scalar weight, let $R$ be either $r$ or $2r$.
If $|\theta_j-x_j|+2R\le\mathsf t_j$, the shift bound gives
$|\theta_j-x_j+\zeta_j(v)|<\mathsf t_j$ for $|v|<R$.
This numerical tube inclusion uses the interior weight at its
own level; the actual complex formation still uses its stated
field and current domains.

For an ordinary source of size less than $2r$, the coefficient
derivative bound on that disc also gives, when $h_y^{(j)}(p)\ne0$,
\[
 \begin{split}
 |b_j(W(p_y))-b_j(W(p))|
 &\le C d_j M h_j\operatorname{Lip}W\\
 &\le C d_jL^{-(n-j)}\|w\|_{\flat,K-n},\qquad j\le n.
 \end{split}
\]
This freezes only the source value in the coefficient. It does
not differentiate the stage profiles.

When $x_j(p)\ge x_{\min}>0$, the modulus of the shifted Wilson weight
at a real field is bounded by
\begin{equation}\label{si:eq-real-wilson-source}
 |e^{-A(x_j-\zeta_j,U)}|
  \le e^{-(1-4\|W\|_\infty/(3x_{\min}))A(x_j,U)}.
\end{equation}
This follows plaquette by plaquette from $a_p\ge0$. It keeps a fixed
fraction of the real Wilson gain when the source radius is fixed
and $x_{\min}$ is sufficiently large. It does not give a bound for
the retained local integrals with their other sourced factors.

\begin{samepage}
\begin{proposition}[A real shift in the same frame]
\label{si:prop-same-frame-source}
Suppose a real prefix has weight-tube room $2r$, and fix a real
$|v|<r$. Use the same frozen tests and reference laws for the
shifted presentation. Assume the complete matching hypotheses
of \cite[Proposition~7.1]{RT} at each formed step: hereditary
physical field--current reads, exact common centered regular
formation with the omitted earlier remainders read on their
actual pairs, the common vacuum convention and cofinal physical
coefficient readout. In particular the composed incoming exponent
and omitted remainder totals are holomorphic on a common connected
complex fluctuation domain containing zero and the real cut.
A nonempty full-dimensional real open patch maps their reads into
the hereditary windows. Use the common positive continuous
density formula in a reference measure with full support there.
Assume a bound $r/3$ on each full increment beyond its shared
one-loop term. Require initial admissibility at $\beta-v$ and
actual local successor admission in this common frame whenever
a formed next weight lies in the tube. This includes the geometric,
field and retained tests on the cofinal compatible tori.
Then the run at $\beta-v$ in this frame is admitted through the
same prefix. Its interior weights and
coefficients satisfy
\begin{equation}\label{si:eq-same-frame-source}
 x_j(\beta-v)=x_j(\beta)-\zeta_j(v),\qquad
 b_j(\beta-v)=b_j(v).
\end{equation}
If $E_{(j)}$ is the once-counted vacuum at the stop and the source
energies are those of \eqref{si:eq-local-energy}, then
\begin{equation}\label{si:eq-real-source-vacuum}
 E_{(j)}[\beta-v]=E_{(j)}[\beta]+\sum_{i\le j}\sum_X e_{K-i,X}(v).
\end{equation}
\end{proposition}
\end{samepage}
\begin{proof}
At the bare level the two input measures agree. Their actual
positive kernels through any common admitted prefix agree, so their
history terms, including the all-small-field term, agree. Inductively
suppose \eqref{si:eq-same-frame-source} holds through level $i-1$.
Formation at the admitted input for step $i$ can be applied before
claiming admission of its output. The possible discrepancy in the
next weight is bounded by $(4/3)|v|+2r/3<2r$. The second term is the
sum of the two errors beyond the common one-loop coefficient. Thus
the new entry lies in the tube. The next-formation clause of
\cite[Proposition~7.1]{RT}, using the common reads and complex
continuation just stipulated, identifies the two coefficients
before closure at this new weight.
The update then gives the weight equality at $i$. Invoke the
stated successor admission to close the presentation at this
entry before using its density. The full
all-small-field densities are continuous on the common open
identity window and are equal as measures there in a common
reference measure with full support. They therefore agree at the
identity. The centered activities and actions vanish there, and
the fixed gauge and coordinate factors coincide. The value of the
exponent at that point is the once-counted vacuum, including its
local source differences in the sourced presentation. This proves
\eqref{si:eq-real-source-vacuum}. No termwise equality of all
localized activities in the two presentations is needed.
\end{proof}

At a prescribed stop, differentiation along this common admitted
real interval gives $\partial_\beta x_j(\beta)=\zeta_j'(0)$,
with the reference frame held fixed. The sum in
\eqref{si:eq-source-slopes} bounds its departure from one at that
stop, rather than only at the end of a longer trajectory.


\section{Ordinary formation with source-dependent members}
\label{si:sec-member-formation}

Fix $q>p_0\ge0$ and write
\[
 \alpha_j=\theta_j^{-1/2}(\log\theta_j)^q,
 \qquad \delta_j=\theta_j^{-1/2}(\log\theta_j)^{p_0},
 \qquad s_j=\delta_j/\alpha_j.
\]

The source is inserted into the Wilson potential at the fixed
reference. The local Gaussian precision, its normalizer and the
retained tests are the same in the sourced and unsourced integrals.

A localization domain $X$ is a connected union of cubes at the
specified level. Write $d_j(X)$ for the shortest cube-edge tree
within $X$ that visits its cubes. A member $F(X,P,z)$ has a
pointing $z$ at a site of $X$ and a pair argument $P=(U,J)$,
where complex bond matrices of $U$ lie in $\operatorname{SL}_2(\mathbb C)$.
The current $J$ is a bond one-form valued in the complexification
of the traceless Hermitian colour space, with the factor-$i$
convention of \cite{EM}. The pair ranges over the actual
complex core of that member, including its declared restrictions
and companion, when used. For a pointed family the exponential
norm at rate $\gamma$ is
\[
    \|F\|_{j,\gamma}
       =\sup_{X,P,z} e^{\gamma d_j(X)}|F(X,P,z)|.
\]
For an unpointed family omit $z$. For a potential that also has
a fluctuation argument take the supremum over its actual allowed
fluctuation coordinates. A rooted sum is a sum over the domains
through a fixed cube or assigned pointing. We only use such sums
with a strict decay reserve above their counting rate.

An ordinary window is a local formation chart at the fixed real
reference, on which all vertices use their specified complex
cores. Its complete determining read comprises every coordinate,
restriction and guard consulted by the operation, including the
reads of its old vertices. The unmarked regular family is formed
without the marked old remainders. Subtracting it from the full
regular family leaves the marked part, with its inherited reads.
Retained operations use their own stored variables and local
measures. These are the conventions for the actual regular and
history operations of \cite{EM,RT}.
At each step the regular output is formed first. Its constant-source
companion selects the coefficient to subtract at each pointing;
the resulting weights are then fixed before closing the new retained
and remainder terms.

We use the following local analytic inputs. The localized
interaction logarithm is formed on the actual Gaussian and
switch domains, with a fixed spare slot for a potential of norm
less than $w_*>0$. Its connected series converges normally and
has a uniform exponential norm there. The assigned bond
pointing has its own local estimate on the same complete reads.
The window bounds for inherited members use the actual complex
potential and current domains of those members. A calculation
at complete physical pairs does not admit independent currents.

\begin{samepage}
\begin{lemma}[The fixed-reference source vertex]\label{si:lem-source-vertex}
Let $\eta$ be the bare mesh in the units of a step and $0<\delta\le\alpha$. Suppose the
actual local pair substitution and its decoupled versions have
plaquette matrices within $C\alpha\eta^2$ of the identity and the
fluctuation changes a plaquette matrix by at most $C\delta\eta^2$.
Suppose the assigned physical bracket at an independent pair has
the same admitted bounds. Then the Wilson source vertex after
subtracting its background value has local norm at most
$C M^4\alpha\delta\|W\|_\infty$. In the ordinary local expansion,
the unpointed derivative in an inserted potential of norm $\epsilon$
has exponential norm at most $C\epsilon$ at the output rate of the connected
series. At a complete physical pair its sum is the expectation of
the inserted potential in the normalized tested regular integral.
The Gaussian normalizer and the scalar ratio at the fixed real
weight have zero source derivative.
\end{lemma}
\end{samepage}
\begin{proof}
For $Q\in\operatorname{SL}_2(\mathbb C)$,
$1-\tfrac12\operatorname{Tr}Q=\tfrac12\det(I-Q)$.
The difference of determinants of two $2\times2$ matrices is
bounded by a constant times their difference in norm times the
sum of their norms. Apply this to the two plaquette differences
from $I$. The displacement and $\delta\le\alpha$ bound the second matrix by $C\alpha\eta^2$ even if the radius is given at the first. The change of one Wilson density is at most
$C\alpha\delta\eta^4$. There are at most $C M^4\eta^{-4}$
plaquettes in an assigned local cube and its fixed collar. This
proves the vertex bound, including complex substitutions. At an
independent pair the assertion uses the admitted local bracket in
place of its physical specialization.

For a potential direction $v$ with norm $\epsilon>0$, insert $t v$
in the actual localized expansion on $|t|<w_*/\epsilon$.
The derivative of a product of selected Mayer factors is the sum
with one factor $e^{V(Y)}-1$ replaced by $v(Y)e^{V(Y)}$.
The corresponding tests and filing domain remain in that term,
including when $V(Y)=0$. The fixed Gaussian domination and normal
convergence justify differentiation term by term. If the bound
of the local logarithm on this disc is $B e^{-\gamma d(X)}$,
Cauchy on $|t|=w_*/(2\epsilon)$ gives
$2 B\epsilon e^{-\gamma d(X)}/w_*$. A zero direction contributes
zero. The same circles bound the norm Taylor remainder by
$C B\epsilon^2/w_*^2$ for $\epsilon<w_*/2$. Thus the bounded
linear first jet also differentiates an actual source curve of
potentials in this norm. The auxiliary $t$-disc need not be a
physical source disc. The renewed pointing uses the stated
rooted pointing estimate, which is separate from the unpointed
Cauchy bound.

At a complete physical pair the exponential of the sum of the
local logarithms is the tested Gaussian integral $\Xi(t)$.
The connected expansion provides its nonzero branch there.
Differentiating the integral and this identity divides the
inserted integral by $\Xi(0)$, giving the normalized expectation.
All Gaussian scalar terms are fixed and therefore differentiate
to zero. No division by a full history mass is involved.
\end{proof}

\begin{samepage}
\begin{lemma}[Rooted localization of the Wilson insertion]
\label{si:lem-source-switches}
For each plaquette $p$ let $C_p$ be its assigned, bounded
face-connected base cubes. Use the
actual hereditary local response brackets on every set $Y\supset C_p$.
In the response formula let $s_D$ be the fine cube switches, and
let $\nu$ multiply the fluctuation response. Assume the plaquette
source evaluation $F_{p,Y}(s,\nu)$ is jointly holomorphic on
$|s_D|<e^\gamma$, $|\nu|<\rho$, where $\gamma\ge1$ and
$\rho\ge c\alpha/\delta>1$. Assume the complete pair substitutions
are admitted there with
$|F_{p,Y}(s,\nu)|\le C\alpha^2\eta^4|W(p)|$. The base switches are
one. The response has the usual restriction property when switches
outside $Y$ are zero. A differenced switch in a component disjoint
from the base gives zero. Then the filed scalar-source potential
has exponential supremum norm at most $CM^4\alpha\delta\|W\|_\infty$ at any fixed
rate below $\gamma-1$. At a physical pair its sum is the actual
Wilson difference between the fluctuating and reference fields.
\end{lemma}
\end{samepage}
\begin{proof}
Put $f(s)=F_{p,Y}(s,1)-F_{p,Y}(s,0)$. Schwarz in $\nu$ bounds
$|f(s)|$ by $2C\rho^{-1}\alpha^2\eta^4|W(p)|$. For the term filed
on $Y$ integrate $\prod_{D\in Y\setminus C_p}\partial_{s_D}f$
over the real unit switch box, keeping the base fixed. Cauchy at
each real switch value has radius approaching $e^\gamma-1$.
The absolute bound on this bracket is
\[
 C\alpha\delta\eta^4|W(p)|
                    e^{-(\gamma-1)|Y\setminus C_p|}.
\]
Each switch is differentiated once, so no factorial is present.
By the response rooting property every nonzero component meets
the base. Connect the base by a bounded tree and add each other
cube through its face-adjacent predecessor, once per cube. This
gives $d(Y)\le |Y\setminus C_p|+O(1)$.
The number of base plaquettes in a filed union is at most
$CM^4\eta^{-4}(1+d(Y))$. Absorb this last factor in the rate
reserve to get the claimed norm. Repeated use of the fundamental
theorem of calculus in the switches expresses the all-one
response as the sum of these brackets, including the base-only
term. Hereditary restriction identifies these summands with the
actual local rule. At a physical pair the endpoints $\nu=1,0$
are the fluctuating and reference Wilson terms, respectively.
This recovers their exact difference. Independent current pairs
only use the admitted brackets and need no global physical
substitution.
\end{proof}

\begin{samepage}
\begin{lemma}[First variations at fixed reference]
\label{si:lem-source-first-variations}
Let $d\lambda$ be a fixed positive sigma-finite measure containing
the Gaussian and Borel test factors of a local leg. On a complex
disc about $t=0$, let $V_t,L_t$ be jointly measurable in parameter
and field, and holomorphic in $t$ at almost every field. Suppose
$L_t$ is locally uniformly bounded in the field supremum norm
and $e^{V_t}$ has a common integrable majorant on each compact
parameter disc. Suppose the denominator below is nonzero near
zero, and set
\[
 \langle F\rangle_t=
 \frac{\int F e^{V_t}\,d\lambda}{\int e^{V_t}\,d\lambda}.
\]
Dots denote derivatives at zero. Then all three terms on the
right below are absolutely defined, and
\begin{equation}\label{si:eq-source-leg-derivative}
 \left.\partial_t\langle L_t\rangle_t\right|_0
   =\langle\dot L\rangle_0+\langle L_0\dot V\rangle_0
                   -\langle L_0\rangle_0\langle\dot V\rangle_0.
\end{equation}
For the fixed pointed subtraction in \eqref{si:eq-member-subtraction},
let $G_t$ be differentiable in the member norm and let $b$ be the
fixed bounded linear readout. Its derivative uses its own
coefficient $b(\dot G)$. If $b^t$ is differentiable in bounded
operator norm and the classical action $A^t$ is differentiable
in a common field chart, then
\begin{equation}\label{si:eq-moving-readout}
 \left.\partial_t\bigl[b^t(G_t)A^t\bigr]\right|_0
 =\bigl[\dot b(G_0)+b^0(\dot G)\bigr]A^0+b^0(G_0)\dot A
\end{equation}
on any common differentiable field chart.
\end{lemma}
\end{samepage}
\begin{proof}
Cauchy's formula bounds $\dot L$ in the field supremum norm.
The same formula applied to $e^{V_t}$, with its common integrable
majorant, gives an integrable bound for $\dot V e^{V_0}$.
Multiplying the latter by bounded $L_0$ proves the separate
absolute integrability claimed in the statement. Cauchy and
Fubini also differentiate both original integrals. Differentiating numerator and
denominator gives \eqref{si:eq-source-leg-derivative} by the
quotient rule. Thus the bound for the renewed bond pointing must
include the direct insertion derivative and the connected term
with the potential derivative on its actual read. The fixed
readout and vacuum evaluation are bounded linear operations
in the pointed member norm, so differentiation commutes with
them. Differentiating the two moving operations by linearity
and the product rule gives \eqref{si:eq-moving-readout}. In the
frozen source presentation both additional reference derivatives
vanish.
\end{proof}

Write $X^+$ for the complete enlarged source read of a domain. The
geometric input is the native one-layer enlargement with one
fixed additional bare-plaquette reach for barycentres. All local
operations are evaluated on complete determining sets. The actual
filing rule puts a fine vertex in a regular output domain only
when the coarse hull of its complete one-layer read lies in that
domain. Boundary and retained terms use their own full collars.
At the
filing of an output $(X,z)$ the enlarged reads of its vertices,
including earlier complete reads, lie in $X^+$. The same condition
holds for the read of an inherited boundary summand. A retained
component uses its own collared carrier and its fixed stored
variables. When it is included in a new filed output, that full
read is part of the output carrier. These are the geometric
admission conditions on the formation. They also give
\begin{equation}\label{si:eq-source-diameter}
 \sup_{p\in X^+}d(x(p),x(p_z))\le C M h_j(1+d_j(X)).
\end{equation}
The path of a cube tree has length at most its number of edges
times $M h_j$, and the endpoint and collar errors are at most a
fixed multiple of $M h_j$.

\begin{samepage}
\begin{lemma}[Comparison on a local data line]\label{si:lem-source-data-line}
Suppose a filed output reads only its complete input stock on
$X^+$. Compare a variable source $W$ there with a constant $v$,
and put
\[
 m=\max\{\sup_{p\in X^+}|W(p)-v|,\ c_M h_j\operatorname{Lip}W\}.
\]
On the independent local data box suppose the line from the
constant data to the sourced data, with the common reference
fixed, is admitted for $|t|m<r$. When that disc contains $1$, the
line has the actual sourced output at $t=1$. Suppose the
normalizer-subtracted output is holomorphic on that disc and is
bounded on the line and at both actual endpoints by
$B e^{-\gamma d_j(X)}$. Its endpoint
difference is at most
\begin{equation}\label{si:eq-data-line-slope}
 2B\min\{1,m/r\} e^{-\gamma d_j(X)}.
\end{equation}
With $v=W(p_z)$, $\gamma\ge\kappa+1$, and $j\le n$, the right
side is bounded by
\begin{equation}\label{si:eq-source-gradient}
 C B r^{-1}L^{-(n-j)}\|w\|_{\flat,K-n}\,e^{-\kappa d_j(X)}.
\end{equation}
\end{lemma}
\end{samepage}
\begin{proof}
For $0<m<r$ subtract the value at $t=0$ and apply Schwarz on
$|t|<r/m$. For $m\ge r$ use the two endpoint bounds. At $m=0$ the
local input reads coincide. With the anchor value as constant,
\eqref{si:eq-source-diameter} bounds $m$ by
$C(1+d_j(X))c_M h_j\operatorname{Lip}W$. Absorb the diameter using
the spare unit of decay and use $h_j/h_n=L^{-(n-j)}$ and the native
source norm at depth $K-n$. This proves \eqref{si:eq-source-gradient}.
The line is in the independent data box. Interpolation in the
physical source class alone would not give this estimate up to its
boundary.
\end{proof}

\begin{samepage}
\begin{proposition}[Formation and old-origin summation]
\label{si:prop-member-formation}
Use the actual regular formation, component operations and
restrictions just described. At level $n$ assume the ordinary
window estimates on their admitted field domains. Let
$\rho_{jn}\ge1$ bound the ratio of the radius required by the
output window to the available radius of its old member;
the value one is used when that ratio is smaller. A single
normalized regular origin at level $j$, with its own vacuum and
action subtraction, contributes at most
$C G_*\rho_{jn}^3 L^{-(4+\alpha_H)(n-j)}$ to the window, where
$\alpha_H>0$. A vacuum-subtracted defect or remainder with norm
$B_j$ contributes at most $C B_j\rho_{jn}^2L^{-4(n-j)}$.
The bounds include the sums over domains through that origin and
$\rho_{jn}^2\le C(1+n-j)$. Assume the complete filed data boxes
admit their constant and variable endpoints, and that their local
lines satisfy the hypotheses of \cref{si:lem-source-data-line},
including holomorphy on the full independent disc. Assume the
actual normalizer-subtracted pointed output is bounded by
$\tfrac12 e^{-(\kappa+1)d_j(X)}$ on those lines and endpoints, and
remainders and boundary outputs have their respective ordinary
rooted bounds there. On each source family the successive outputs
are holomorphic, are real for real sources and fields, and depend
on the source only through their complete bare plaquette reads.
The difference from the selected constant-source member is a
vacuum-subtractable defect of norm at most
$C r^{-1}\|w\|_{\flat,K-n}L^{-(n-j)}$ at age $n-j$.
If the vacuum-subtracted remainder norms are at most
$C_R\theta_j^{-p}$, $p>2$, and the lower drift in
\cref{si:prop-source-profiles} holds, the total ordinary input from these
three families is bounded by
\begin{equation}\label{si:eq-member-sum}
 C G_*+ C r^{-1}\|w\|_{\flat,K-n}+C C_R\theta_n^{2-p}.
\end{equation}
The bounds hold for the sourced filing with a different selected
member at each base point.
\end{proposition}
\end{samepage}
\begin{proof}
At the bare stage the source enters only as the values of $W$ at
plaquette points. At a subsequent stage first form the actual
regular output with the ordinary rule. The complete input
reads of a connected term are inside its filing stock. The
source values of an earlier read are there too by the geometric
inclusions. Products, fixed coordinate substitutions, integration
and normally convergent rooted sums therefore give holomorphy
and dependence on these values alone. Conjugation commutes with
these operations and fixes their real reference branches. The
same reasoning on each complete retained carrier gives the
boundary and remainder assertion, using the supplied domination
for those operations. Carried terms are restricted with their
original carriers and branches. Thus equality of two source
representations on the complete read gives equality of the output.

At each new pointing select the whole-lattice member at its
anchor value. It is formed with the same ordinary interior rule.
When $W$ is constant on the complete read there is equality
with this member on the common field window. Otherwise
\cref{si:lem-source-data-line} gives its defect bound. Taking out
the member's own coefficient and reinstating it by
\eqref{si:eq-running-shift} gives the exact exponent, by the
pointwise identity for the weights. Thus at the next step the
actual sourced input consists of normalized selected members,
defects and remainders on their declared domains. Bounds on
these domains and on their local lines are needed for the
construction. Neither holomorphy nor summation by itself supplies
those admissions.

A level-$n$ window contains at most $C L^{4(n-j)}$ origins of
level $j$. In each of the window estimates this cancels the
volume factor. With $a=n-j$, the regular sum is bounded by
$C G_*\sum_{a\ge0}(1+a)^{3/2}L^{-\alpha_H a}$. The defect sum
is bounded by
$C r^{-1}\|w\|_{\flat,K-n}\sum_{a\ge0}(1+a)L^{-a}$.
For remainders, $\theta_j\ge\theta_n+\lambda a-C_0$ and a
large fixed floor bound the sum by a constant times
\[
 C_R\sum_{a\ge0}(1+a)(\theta_n/2+\lambda a)^{-p}
 \le C C_R\theta_n^{2-p}.
\]
The last inequality follows by comparison with the integral
of $(\theta_n/2+\lambda t)^{1-p}$, which is finite for $p>2$.
Every constant used the supremum of the member norms and is
independent of the sitewise selections. This proves
\eqref{si:eq-member-sum}.
\end{proof}

\paragraph{Admission of the comparison lines.}
For full ordinary source radius $2r$, use the constant companions
on $|v|<4r$. The argument of \cref{si:prop-source-profiles} on this
larger disc uses $6r$ of half-product room, with the marked estimate
on the corresponding actual companion disc. It gives the same
summable bound on the radial slopes. Cauchy therefore gives a
Lipschitz constant $2d_i$ for $b_i$ on $|v|<2r$. This reserve is
chosen in the supplied local formation before the inverse-coupling
floor.

Here is a sufficient check of the independent line hypotheses in
\cref{si:prop-member-formation}. Suppose the analytic ordinary rule admits
site-source and running-weight slots of radius $4r$. On the native
carriers let $B_i$ bound the normalizer-subtracted regular,
boundary and remainder slots, with their actual rates and
remainder small factors. Assume the enlarged box admits, about each constant-source
centre, perturbations of norm $2B_i$ in these slots, measured
with their native rates and small factors. In a box about zero,
$3B_i$ of absolute room suffices. Give the independently
selected members the same comparison reserve. The
ordinary formation estimates on this box must bound the output
regular legs in the following exponential norms. With the decay rate
$\kappa+1$, the interaction leg costs $Cs_k$ times the sum of the
old-origin input norms and $Cs_k\alpha_k^2$ times the Wilson shift
norm. The intrinsic, coordinate and cutoff legs together cost
$C M^4\theta_k^{-1/2}\mathfrak b_k^3+C\theta_k^{-1/4}$, where
$\mathfrak b_k$ is at most a fixed power of $\log\theta_k$. The free
quadratic correction costs $C\mathsf t_k/\theta_k$ on the same rate,
where $\mathsf t_k$ is the real tube radius. These are estimates
for the actual formed legs on their local cores. The two
reference Gaussian contributions being subtracted are common
to the line. For a retained component use the absolute excess
recursion and domination on this enlarged box, as in the
retained-integral hypotheses below.

Fix a complete output stock and a constant comparison value
$v$. On the stock interpolate the site values, selected regular
members and their own coefficients affinely from constant to
variable data. Retain the fixed reference entries elsewhere.
Take $m$ as in \cref{si:lem-source-data-line} at the output scale.
For $|t|m<r$ a site value has modulus at most $3r$. The sum of
the coefficient departures from the zero-source values, including
the interpolated differences, is at most
$(2r+2r)\sum_i d_i\le2r/3$. Thus the running weight shift fits the
$4r$ slot. By induction the input regular and boundary differences
at their respective complete reads are bounded by $2B_i m/r$.
Multiplication by $t$ fits their reserved room. The same holds for
the remainder slots with their own small factors. An old defect
vanishes at the constant endpoint. Its older bound from
\eqref{si:eq-source-gradient}, applied before the final source
norm comparison, is $C r^{-1} h_i\operatorname{Lip}W$. Thus its
line norm at output level $k+1$ is at most $CL^{-(k+1-i)}$.
These checks use only complete reads inside the chosen stock. The
linear restrictions of a boundary sum are applied to its individual
summands. A retained excess is generated by its actual operation
on the inherited lines, including the reference contribution and
all earlier excess terms. This uses domination on the full input
line. Interpolating only a part of an exponent could lose a
cancellation between terms with large constant parts.

The ordinary origin sums are now uniformly bounded by the
geometric, derivative-of-geometric and power sums in
\eqref{si:eq-member-sum}. The displayed leg estimates give a
normalizer-subtracted regular output on the line bounded by
\[
 C\bigl(s_k+M^4\theta_k^{-1/2}\mathfrak b_k^3
                +\theta_k^{-1/4}+\mathsf t_k/\theta_k\bigr)
                     e^{-(\kappa+1)d_{k+1}(X)}.
\]
A sufficiently small relative tube and sufficiently large floor
make the prefactor at most $1/2$. The ordinary bounded output
estimates on the enlarged box give the other $B_{k+1}$ bounds.
Applying \cref{si:lem-source-data-line} recovers the required
slopes for the next stage. At the first level the bare site
interpolation is affine and there are no old regular or retained
entries. This inductive check uses the local estimates on the
actual enlarged box. It deduces source-line bounds from them
rather than extending the physical-source disc by assertion.

\subsection{First source derivatives and the near-window domains}
\label{si:subsec-near-source}

The near-window conclusion below uses a complete nonwrapping
background-plus-fluctuation datum of size at most $C\delta_n$ at
level $n$, including its gauge and remote reads. We specify the
local background input for an origin of level $j\le n$. Put
$\sigma=L^{-(n-j)}$ and $D=1+d_j(X)$. At each fixed pointing $z$
use one affine curvature $F_{j,z}$ and its linear affine pair $L_F$
in the common old window chart. Every domain in the pointed total,
including the tent action, uses the restriction of this same pair.
In the old pair norm the evaluated near pair $P$ satisfies
\begin{equation}\label{si:eq-near-affine-pair}
 \begin{split}
 \|L_F\|_X&\le C D^b\delta_n\sigma^2,\\
 \|P-L_F\|_X&\le C D^b(\delta_n\sigma^{2+\alpha_H}
                       +r_{B,j}^{-1}\delta_n^2\sigma^4),
 \qquad \alpha_j/r_{B,j}\le C,
 \end{split}
\end{equation}
where $\alpha_H>0$ and $r_{B,j}$ is the old independent datum
radius. The settled lower plaquette bound is $C\delta_n\sigma^2$.
The common upper and remote reads have amplitude $C\delta_n$.
Use the actual weighted response on those reads and the full
absolute datum ceiling, with every complex substitution admitted.
The affine norm bound in \eqref{si:eq-near-affine-pair} holds also
on domains omitted from the near chart. Any such omission used
in the affine quadratic cancellation has length at least
$c/\sigma-C$.

Each local old regular and defect member, at its individual
$(X,z)$, is invariant under simultaneous constant adjoint action
on its admitted pair chart. For the complete constant-source
member use the affine coefficient identity at each pointing of
\cite[Lemma~7.1]{EM}, including its scalar lattice covariances.
The identity for complex affine curvatures follows from its real
version by polynomial complexification of the quadratic forms.
Separately from the datum radius, the old regular and defect
members, and their first-source jets used below, are holomorphic
on pair charts containing a ball of radius $c\alpha_j$
about the trivial pair in their declared norm. Their norm bounds
hold on these balls and on all actual evaluated pairs, even when
an evaluated pair lies outside the smaller Taylor ball. The
pointed action is holomorphic on the field projections of these
balls and is admitted on all actual evaluated fields. Its norm is
$O(\alpha_j^2)$ on both sets on its fixed bounded carrier.
The affine readout of a unit curvature has pair norm bounded by
$C(1+d_j(X))^b$ on each old domain. Its transported Hessian
uses the unit coefficient convention of
\eqref{si:eq-member-subtraction}. The field estimates for the
complete window and its independent current are the local inputs
here. The size of a fluctuation alone does not supply them.

\begin{samepage}
\begin{proposition}[First-source response at fixed reference]
\label{si:prop-first-source-response}
Assume the lower drift in \cref{si:prop-source-profiles} and
$\theta_i-\theta_j\le C(1+j-i)$ for $i\le j$. For a source
direction $f$, write $D_j[f]=\partial_t E_j^{tf}|_{t=0}$ for the
pointed regular jet, and let $A_j$ be its operator norm from
$\|f\|_{\sharp,K-j}$ to $\|\cdot\|_{j,\kappa+1}$. The initial
regular list is empty. Use the actual one-step vertex and renewed
pointing estimates on complete admitted reads, with the following
weights. The Wilson direction costs $C\alpha_k^2s_k$ times the
shift norm. The fluctuation difference of an old normalized
regular jet has weight $Cs_k\omega_{k+1-j}$, where $\omega_a\ge0$
and $\sum_{a\ge1}\omega_a<\infty$. An old vacuum-subtracted
defect or admitted aged remainder has weight
$Cs_k(\alpha_k/\alpha_j)^2$. For every remainder of origin $j$
and every source direction $f$, its actual reads are admitted on
the line $|t|\|f\|_{\sharp,K-j}<r$, its value is holomorphic there,
and its norm there is at most
$C_R\theta_j^{-p}$ with $p>3$. This radius uses the physical
scale of that origin, not the unit-scale test norm.
For a sufficiently large floor these estimates imply
\begin{equation}\label{si:eq-first-source-amplitude}
 A_j\le C\alpha_{j-1}^2s_{j-1}\qquad(j\ge1).
\end{equation}
If the near-window inputs of \cref{si:subsec-near-source} hold,
the total of the renormalized near regular and defect jets per
level-$n$ window is at most
$C\delta_n^2s_n\|f\|_{\sharp,K-n}$, where
$\delta_n^2s_n=\theta_n^{-1}(\log\theta_n)^{3p_0-q}$.
Additional winding, marked and
field-independent remote tails retain their own bounds.
\end{proposition}
\end{samepage}
\begin{proof}
Subtract $f(p_z)D_j[1]$ from the jet at $(X,z)$. By
\eqref{si:eq-source-diameter}, the profile $f-f(p_z)$ on the
complete read has size at most
$C h_j(1+d_j(X))\operatorname{Lip}f$. Clip its real and imaginary
parts at that bound outside the read. The clipped profile has
unchanged values on the read and Lipschitz constant at most a
fixed multiple of $\operatorname{Lip}f$. The jet is linear and
local. The clipped sharp norm and the spare decay unit give the
bound $C A_jh_j\operatorname{Lip}f$ for the defect at rate
$\kappa$. At level $k+1$ this is at most
$C A_j L^{-(k+1-j)}\|f\|_{\sharp,K-k-1}$. The derivative of the
selected member's coefficient is its own $b_j'(0)f(p_z)$ by
linearity of the fixed physical readout. Differentiating
\eqref{si:eq-running-shift} bounds the actual Wilson shift
derivative by $(1+\sum_jd_j)\|f\|_\infty$. Cauchy on the
remainder disc bounds the remainder derivative by
$C r^{-1}\theta_j^{-p}\|f\|_{\sharp,K-j}$.

Let $r_{kj}=\alpha_k/\alpha_j$. The one-step bounds give
\begin{equation}\label{si:eq-first-source-recursion}
 A_{k+1}\le C s_k\left[\alpha_k^2+\sum_{j=1}^k\omega_{k+1-j}A_j
   +\sum_{j=1}^k r_{kj}^2 A_j L^{-(k+1-j)}
   +r^{-1}\sum_{j=0}^k r_{kj}^2\theta_j^{-p}\right].
\end{equation}
The norm of $f$ at an earlier level is bounded by its norm at
level $k+1$. The adjacent weights have bounded ratios. For large
$\theta$, both $\alpha(\theta)^2$ and $s(\theta)$ decrease, and
$\theta_j\ge\theta_k-C_0$ for $j\le k$. Thus
$\alpha_{j-1}^2s_{j-1}\le C\alpha_k^2s_k$. With $r_{kj}^2$
inserted the same bound follows from
$r_{kj}^2\alpha_{j-1}^2s_{j-1}
 =\alpha_k^2(\alpha_{j-1}/\alpha_j)^2s_{j-1}$.
The upper drift gives $r_{kj}^2\le C(1+k-j)$. The last sum in
\eqref{si:eq-first-source-recursion} is $O(\theta_k^{2-p})$ as in
\eqref{si:eq-member-sum}. Since $p>3$ it is $O(\alpha_k^2)$ at
fixed $r$. Under the inductive bound
$A_j\le C_*\alpha_{j-1}^2s_{j-1}$ the right side is at most
$\alpha_k^2s_k(C_1+C_2 C_*s_k)$. Take $C_*\ge2C_1$ and make
$C_2s_k\le1/2$ by the floor. Empty sums start the induction.

For the near conclusion, the central axial comb uses at most
$C\sigma^{-1}$ plaquette factors of size $C\delta_n\sigma^2$.
Their product has link logarithm of size $C\delta_n\sigma$.
The cutoff extension costs $C\delta_n$ in the full row norm.
The remote rows have the proportional amplitude $C\delta_n$,
so their weighted first response is bounded by
$C\delta_n e^{-c/\sigma}$ and fits the first error in
\eqref{si:eq-near-affine-pair}. The central half-disc condition
is uniform in age because
\[
 \delta_n\sigma/r_{B,j}
  \le C s_n(1+n-j)^{1/2}L^{-(n-j)}.
\]
The right side is small at the chosen floor. The remaining
absolute core inclusions are those of the standing window data.

The regular member and the defect both have zero linear term
at the trivial pair. The ordinary scalar source preserves their
constant adjoint invariance at each pointing, and the adjoint
representation has no invariant linear functional. In the
complete regular pointed sum the Hessian on $L_F$ is cancelled
by its own tent action at that pointing. The coefficient of a
member of norm $G$ is bounded by $C G\alpha_j^{-2}$. To see this,
apply its Hessian to the unit affine curvature used for the
pointed coefficient. The affine pair on a domain grows at most
polynomially in its length, the pair ball gives the factor
$\alpha_j^{-2}$ by Cauchy, and the domain sum converges by the
strict entropy margin. The unit action Hessian is fixed. Thus
the own tent subtraction has norm at most $C G$, by its
$O(\alpha_j^2)$ bound. This estimate uses the fixed affine readout
on the complete pointed member.

Put $\rho_j=\delta_n/\alpha_j$ and $u=C D^b\rho_j\sigma^2$.
For small $u$, \eqref{si:eq-near-affine-pair} places $P$ in the
half ball of the old pair chart. For a local term of norm $g_X$
including its exponential decay, Cauchy's estimate gives a
Taylor remainder $C g_Xu^3$. The difference between its
quadratic values on $P$ and $L_F$ is at most
\[
 C g_X D^{b'}(\rho_j^2\sigma^{4+\alpha_H}
                         +\rho_j^3\sigma^6).
\]
The two products of the affine field with the errors give these
terms directly. The squared errors satisfy the same bound,
using $u$ small and $\sigma\le1$. For $u$ above a fixed positive
threshold, use the actual core bound and the affine quadratic
bound $C g_Xu^2$. Their sum is bounded by $C g_Xu^3$.
For the vacuum-subtracted defect, the zero linear term and the
same pair bound give $C g_Xu^2$ for small $u$, and the actual
core bound gives it for large $u$.

Apply the affine cancellation at each pointing before multiplying
by $f(p_z)$. If the near part omits domains, the missing affine
Hessians are bounded by the same $C g_X D^{b'}\rho_j^2\sigma^4$.
Their lengths are at least $c/\sigma-C$, so the rooted sum
adds $O(e^{-c'/\sigma})$ and fits the $\sigma^{4+\alpha_H}$ term.
This use of the omitted affine values imposes no estimate on
an actual far evaluated field. All polynomial factors in $D$
are summable at the reserved decay. There are $O(\sigma^{-4})$
old base points in the window. After this count the regular
factors are $\rho_j^2\sigma^{\alpha_H}$ and $\rho_j^3\sigma^2$,
and the defect factor is $\rho_j^2$.

The constant-source direction at a base point is $f(p_z)D_j[1]$,
of norm at most $A_j\|f\|_\infty$. By the proved amplitude,
$A_j/\alpha_j^2\le Cs_n$ and
$\rho_j\le Cs_n(1+n-j)^{1/2}$. With $a=n-j$, the sum of the
regular terms is bounded by
\[
 C\delta_n^2s_n\|f\|_\infty
     \sum_{a\ge0}\bigl(L^{-\alpha_Ha}+s_n(1+a)^{1/2}L^{-2a}\bigr).
\]
The defect sum is at most
$C\delta_n^2s_n\operatorname{Lip}f\sum_{j\le n}h_j$, which is
$C\delta_n^2s_n\|f\|_{\sharp,K-n}$ by the mesh ratio. The two
estimates prove the stated near bound.
\end{proof}


\subsection{Sources in retained operations}
\label{si:subsec-retained-source}

The complex source changes the integrand of a retained operation on
its fixed stored variables. Write $\Phi_\omega(\xi;P)$ for the
reference exponent of a record $\omega$ at zero source, and write
$\Psi_\omega$ for its source excess. The recorded reference measure
$\mu_\omega$, the integration set and the real stored-field cut are
fixed. The exterior argument $P$ ranges over the actual admitted
set over $\xi$. This set includes the prescribed test compatibility
for a drop-out. The Wilson contribution carried by another factor
stays on that factor.

\sineedspace{45\baselineskip}
\begin{samepage}
\begin{lemma}[Source excess and retained mass]
\label{si:lem-retained-source-excess}
Consider a finite nested family of actual retained operations on a
common admitted source or local-data disc. For each record $\omega$
let $n_\omega^0$ count fresh charged sites in their own scales,
and let $n_\omega=n_\omega^0+\sum_{\omega'}n_{\omega'}$, summing
over the swallowed records. Assume the actual source-sensitive
core Wilson term has the form $A(\widehat c\,\upsilon 1_X,U'')$,
with $|\widehat c|\le\bar c$, $0\le\upsilon\le1$, and real
stored background $U''$. On each nonbare layer $j$ where the
cut is nonzero, assume its plaquette defects are bounded by
$C\eta_j^2$, where $\eta_j$ is the bare mesh in those layer units,
with at most $C\eta_j^{-4}$ such plaquettes per fresh site. On
the bare layer use the uniform plaquette bound and bounded
incidence. Omit this core term when the operation has none.

Suppose the actual fresh collar terms $V_X^{\mathrm{col},W}$ and
$V_X^{\mathrm{col},0}$ have the sum of absolute bounds at most
$C(1+\bar c)n_\omega^0$. This includes their constant parts.
The kept responses of old carrier functions use one common map.
On a kept carrier $R$, applying the response to its excess
function $\psi_R$ gives its value at the actual stored pair
$q(\xi)|_R$ plus a correction $\delta_R[\psi_R]$. Assume the
absolute sum of these corrections is at most
$C(1+\bar c)n_\omega^0$ on the complete reads. Each stored pair
lies in the admitted domain of the corresponding old exponent.
For each record use a fixed Borel integration domain with a
positive sigma-finite reference measure. The evaluated integrands
are jointly measurable on products with compact parameter sets.
The evaluated exponents and reference factors, including the
re-evaluations of old exponents, are holomorphic in the given
parameters for almost every stored tuple on that domain.

Then the source excess of the complete record satisfies
\begin{equation}\label{si:eq-retained-source-excess}
                  |\Psi_\omega|\le C(1+\bar c)n_\omega .
\end{equation}
The constant is independent of nesting depth. Suppose the fixed
reserve weight satisfies $v(\omega)\ge c_v n_\omega$ for a fixed
$c_v>0$. Let $\varpi(\omega)\ge0$ be the actual record weights,
and suppose their real weighted mass satisfies
\begin{equation}\label{si:eq-retained-source-mass}
 \sum_\omega \varpi(\omega)\int^*
 e^{c v(\omega)}\sup_P |e^{\Phi_\omega(\xi;P)}|
                         \,d\mu_\omega(\xi)\le M_c
\end{equation}
for a fixed $c$ large enough to pay \eqref{si:eq-retained-source-excess}.
Then the source-dependent operation tested against a fixed bounded
Borel function of the stored variables has absolute value at most
$M_c\|F\|_\infty$. It is holomorphic on the same admitted parameter
disc for a fixed bounded $F$. Any further visible-record price
uses the same restricted mass at a correspondingly larger reserve.
\end{lemma}
\end{samepage}
\begin{proof}
Subtract the complete zero-source exponent on the same stored
tuple. The fixed reference quadratic cancels. In each kept
response put its value at $q(\xi)|_R$ back in the corresponding
old exponent. The corrections remain in the fresh term, so the
excess is
\begin{align*}
 \Psi_\omega&=\Psi_\omega^{\mathrm{new}}
        +\sum_{\omega'}\Psi_{\omega'}(\xi';q(\xi)),\\
 \Psi_\omega^{\mathrm{new}}
   &=A(\widehat c\,\upsilon 1_X,U'')
     +V_X^{\mathrm{col},W}-V_X^{\mathrm{col},0}
     +\sum_R\delta_R[\psi_R].
\end{align*}
This decomposition uses the complete older exponent at its read.
It does not combine independently chosen pieces of old exponents.
On a nonbare layer the identity
$a_p(U'')=\tfrac12\|U''_p-I\|_{\mathrm{op}}^2$ gives a bound
$C\eta_j^4$ per plaquette. Multiplying by the incidence bound
cancels the mesh factor and costs at most $C\bar c$ per fresh
site. The bare layer gives the same site bound using $a_p\le2$.
Add the two supplied fresh bounds to obtain
$|\Psi_\omega^{\mathrm{new}}|\le C(1+\bar c)n_\omega^0$.
At a leaf this already proves \eqref{si:eq-retained-source-excess}.
For other records add the older estimates at their own counts.
By additivity the new constant is the maximum of the fresh and
old constants. This proves the bound by induction.

For each $\xi$ and each admitted exterior argument, the modulus
of the source multiplier is at most $e^{c v(\omega)}$. Use this
bound before taking the exterior supremum. Equation
\eqref{si:eq-retained-source-mass} then gives an integrable
majorant, including the nonnegative record weights $\varpi$.
Integrating after multiplication by $F$ gives the stated bound.
The same majorant on compact subdiscs permits integration of
the holomorphic integrand around a triangle. Fubini gives zero,
so Morera's theorem gives holomorphy in each parameter. Local
boundedness then gives joint holomorphy on a finite parameter
product. The upper integral is used only to bound the supremum
in \eqref{si:eq-retained-source-mass}. Choose measurable majorants
with total integral arbitrarily close to that bound. The evaluated
integrands themselves are measurable by hypothesis. Restricting records first and
raising $c$ pays a further bounded price per recorded site with
the same argument.
\end{proof}

The absolute collar bounds in this lemma cannot be replaced by
oscillation bounds in the exterior pair. Oscillation does control
the kept-response corrections after their common value at the
stored pair has been subtracted. If the actual normalized carrier
oscillation at the next operation is at most $b_*+\epsilon s$
when the inherited bound is $s$, choose $c_*\ge 2b_*$ and keep
the formation domain $s\le4c_*$. For $8\epsilon\le1$ the new
bound is at most $c_*$. In the ordinary local estimate the
coefficient $\epsilon$ decreases as a positive inverse power of
$\log\theta_j$ at fixed other constants. This is a separate
closure estimate on the common enlarged data domain. It gives
no bound on an omitted field-independent collar value.


\subsection{A common majorant and the nonvanishing boundary}
\label{si:subsec-retained-boundaries}

The retained mass hypothesis uses a common reference coordinate
system before taking a supremum over exterior data. The following
calculation supplies one sufficient envelope when such coordinates
are available. It does not construct those coordinates.

\begin{samepage}
\begin{lemma}[A common Gaussian majorant]
\label{si:lem-common-gaussian}
Let $b\in\mathbb R^m$, where $m$ is a nonnegative integer,
$N\ge0$ and $m\le c_0N$. Let $\sigma$ index
real symmetric matrices $Q_\sigma\ge\gamma I$, where $\gamma>0$.
Suppose $c_0,C_0>0$, $|c_\sigma|\le C_0^m$, $S_\sigma$ is Borel,
and the functions $V_\sigma$ and $H_\sigma(z,\cdot)$ are Borel.
For fixed real constants $B_0,\delta$, assume
\[
 \operatorname{Re}V_\sigma(b)\le B_0N+\gamma|b|^2/4,
 \qquad \operatorname{Re}H_\sigma(z,b)\le\delta
\]
on the parameter set being considered. Then the densities
$1_{S_\sigma}c_\sigma
 e^{-b^{\mathsf T}Q_\sigma b/2+V_\sigma(b)+H_\sigma(z,b)}$
have the common Borel absolute majorant
\[
 C_0^m e^{B_0N+\delta}e^{-\gamma|b|^2/4}.
\]
Its integral is at most $e^{CN+\delta}$, with $C$ depending only
on $c_0,C_0,B_0,\gamma$. The same bound controls the upper integral
of the supremum of the absolute densities.
\end{lemma}
\end{samepage}
\begin{proof}
The lower quadratic bound leaves $-\gamma|b|^2/4$ after paying
the quadratic part of the bound on $V_\sigma$. Taking absolute
values gives the displayed majorant pointwise. Its integral is
\[
 C_0^m e^{B_0N+\delta}(4\pi/\gamma)^{m/2}.
\]
The bound on $m$
absorbs its two dimension factors in $e^{CN}$. The majorant is
Borel even if an uncountable supremum of the densities is not.
\end{proof}

Parameter-dependent changes of variables with separate bounds do
not supply this common envelope. Nor does holomorphic dependence
on a background parameter by itself ensure a nonzero integral.

\begin{samepage}
\begin{lemma}[A phase can cancel a tested integral]
\label{si:lem-phase-zero}
Let $w_0$ be integrable on a measure space, with
$W_0=\int w_0\,d\mu\ne0$. Let $x\ge0$ be measurable, and suppose
the pushforward of $|w_0|\,d\mu$ under $x$ is atomless. There
are $c>0$ and $\lambda\in[-\pi,\pi]$ for which
\[
            \int w_0 e^{i\lambda 1_{\{x>c\}}}\,d\mu=0.
\]
\end{lemma}
\end{samepage}
\begin{proof}
Put $A(c)=\int_{x\le c}w_0\,d\mu$ and $B(c)=W_0-A(c)$.
Atomlessness and integrability make these functions continuous,
with $A(0)=0$ and $A(c)\to W_0$. Thus $|A(c)|-|B(c)|$ changes
from a negative to a positive value at some finite $c>0$.
At a zero, $|A(c)|=|B(c)|>0$; otherwise $W_0=0$.
Choose $e^{i\lambda}=-A(c)/B(c)$ and split the integral at $c$.
\end{proof}

The phase in this example can be independent of every background
parameter, hence holomorphic in those parameters. If $x$ is
gauge invariant, the phase is gauge invariant too. Bounds on
admitted switch and displacement domains, or the explicit local
nonvanishing estimates used below, supply additional information
that this example lacks. A logarithm of a retained or full
partition integral must use such information.


\section{Rooted counting for source reads}
\label{si:sec-source-count}

At spacing $a_s=L^{-s}$ the localization cubes have side $Ma_s$,
with the fixed $M$ of \cref{si:sec-setting}. Take tori of side $\ell=2L^m$ with
$m\ge m_0$, so the localization grids at every retained level tile
the torus. A nonempty domain $X$ is indexed by its set of cubes.
Let $m_s(X)$ be the minimum number of grid edges in a tree
meeting every cube of $X$, allowing a one-vertex tree of length
zero. A full source read $X^+$ lies within a fixed number of
layers of these cubes. The enlargement includes any reads of
plaquette centres.

\begin{samepage}
\begin{lemma}[Rooted domain counts]\label{si:lem-source-count}
For a cube $Q$ and an integer $b\ge0$, the number of domains
$X$ containing $Q$ with $m_s(X)=b$ is at most
$16\,8^{2b}2^{16(b+1)}$. In particular, if $\kappa$ exceeds the
rooted entropy exponent $\kappa_D$, then for some $C,c>0$,
\begin{equation}\label{si:eq-source-count}
 \begin{aligned}
 \sum_{\substack{X^+\cap B_u\ne\varnothing\\m_s(X)\ge t}}
     e^{-\kappa m_s(X)}
   &\le C(1+u/a_s)^4 e^{-ct},\\
 \sum_X e^{-\kappa m_s(X)}
   &\le C\ell^4 a_s^{-4},
 \qquad
 \operatorname{diam}X^+\le C a_s(1+m_s(X)).
 \end{aligned}
\end{equation}
Here $B_u$ is any ball of radius $u$, $t\ge0$, and the distances
are those of the periodic quotient. The same estimates hold for
indexed families with a total rooted count
$C e^{\kappa_D b}$, including their multiplicity.
\end{lemma}
\end{samepage}

\begin{proof}
A grid edge meeting a closed cube has an endpoint in the cube.
Thus a minimizing tree with $b$ edges has a vertex in every
cube of $X$. Start at a corner of $Q$ in the tree, and walk over
each tree edge twice. There are at most $16$ starting corners,
$8^{2b}$ possible walks, and $16(b+1)$ cubes incident to their
vertices. The choice of a subset specifies $X$. On a quotient
there can be fewer distinct cubes and there are still at most
eight oriented steps at each vertex. The same count applies.

If $X^+$ meets the ball, some cube of $X$ is within $CMa_s$ of
it. The number of possible such root cubes is at most
$C(1+u/a_s)^4$. At each root the sum over lengths at least $t$
is geometric, by $\kappa>\kappa_D$. Summing instead over all
root cubes gives the global count. Two points in $X^+$ are
within $CMa_s$ of vertices of a minimizing tree, whose distance
along the tree is at most $Ma_s m_s(X)$. This proves the
diameter bound.
\end{proof}


\section{Local source energies and the common normalization}
\label{si:sec-local-energies}

For each level $j$ and localization domain $X$, write $E^{(j),W}(X,P,z)$
for the newly formed regular member at its pointing $z$, before
its vacuum and classical-action subtraction. Write $R^{(j),W}(X,P)$
for the completed new remainder before its vacuum subtraction.
Use $\Delta_W$ for the difference
from the same member at zero source. Define
\begin{equation}\label{si:eq-local-energy}
 e_{K-j,X}(W)=\sum_{z\in X\cap\Lambda_j}
                \Delta_W E^{(j)}(X,P_\circ,z)
                      +\Delta_W R^{(j)}(X,P_\circ),\qquad j\ge1.
\end{equation}
Here $\Lambda_j$ is the lattice at level $j$. At $j=0$ put $e_{K,X}=0$.
In the source normalizer take the point term $U(W)=0$.

\begin{samepage}
\begin{proposition}[Locality, entropy and the vacuum difference]
\label{si:prop-local-energies}
Use the sourced formation of \cref{si:prop-member-formation}
on the ordinary ball of a prefix, with its complete reads. Assume
that the regular outputs at $P_\circ$, after subtraction of
their common frozen-reference term $N^{(j)}$, have uniform bound
$C e^{-(\kappa+1)d_j(X)}$. Here $N^{(j)}$ is the actual regular
normalizer and assigned fixed-real-weight ratio term, identical
as a function for source $W$ and source zero. Assume the corresponding
remainders have bound $C\theta_j^{-p}e^{-\kappa d_j(X)}$ with $p>0$.
Index a domain by its union of localization cubes, once per union
at that level. Let $\kappa$ exceed the rooted entropy rate for
this indexing. Then the local energies in \eqref{si:eq-local-energy}
are holomorphic, real for real sources and zero at zero source.
They depend on the bare plaquette values in $X^+$ alone and satisfy
\begin{equation}\label{si:eq-energy-decay}
 |e_{K-j,X}(W)|\le C r^{-1}\|w\|_{\flat,K-n}\,e^{-\kappa d_j(X)}
                      \qquad(1\le j\le n)
\end{equation}
on the prefix stopped at $n$. The bound holds also with $m(X)$,
the shortest cube-edge tree length allowed to leave $X$, in
place of $d_j(X)$.

Use the once-counted vacuum convention of \cite{RT}. For this
prefix let $V=E_{(n)}$ be the unsourced vacuum and
$\mathcal E(W)=\sum_{j\le n,X}e_{K-j,X}(W)$. The all-small-field
scalar is $V+\mathcal E(W)$. In any history whose scalar shares are
obtained by the actual filing, let $F_h(0)$ be its unsourced
scalar exponent and let $\mathcal E_h(W)$ be its retained source
share. With this normalization the scalar exponent is
\begin{equation}\label{si:eq-vacuum-share}
  F_h(0)-V-\{\mathcal E(W)-\mathcal E_h(W)\}.
\end{equation}
Suppose each missing pointing and each missing remainder at
level $j$ is incident by its own read to the bounded-layer
neighbourhood of a paid cube of that label. The source difference
in braces costs at most $C\sum_{j\le n}\# Z_j^+$ on the ordinary
ball, where $Z_j^+$ denotes the paid cube set with its fixed collar.
\end{proposition}
\end{samepage}
\begin{proof}
In the once-counted scalar the Gaussian, gauge and Haar factors
of the frozen reference do not vary with the source. All Wilson
terms vanish at the trivial pair. Each regular trivial value is
assigned at its pointing, and the new remainder value is counted
once after closure. Taking their differences gives
\eqref{si:eq-local-energy} and the asserted all-small-field scalar.
The same subtraction in an individual history gives its own
source share. Subtracting the all-small-field scalar then gives
\eqref{si:eq-vacuum-share} with the displayed sign.

Local formation gives holomorphy, reality and the specified reads.
The term $N^{(j)}$ stays fixed when the source changes. At
$P_\circ$ this gives
\[
 \Delta_W E^{(j)}=\Delta_W(E^{(j)}-N^{(j)}).
\]
The remaining trivial value is retained in the vacuum readout.
For the bound, use the holomorphic source line
$tW$ up to the ordinary source radius. Subtract the value at
$t=0$ and apply Schwarz to the uniform bound. Each regular
difference then costs at most
$C r^{-1}\|w\|_{\flat,K-n}e^{-(\kappa+1)d_j(X)}$, and the
remainder difference costs at most this source factor times
$C\theta_j^{-p}e^{-\kappa d_j(X)}$. The zero source gives zero
directly. A tree with $d_j(X)$ edges meets at most $C(1+d_j(X))$
cubes and therefore at most $C M^4(1+d_j(X))$ pointings. The
spare unit of rate absorbs this factor. Since $m(X)\le d_j(X)$,
\eqref{si:eq-energy-decay} also holds at the shorter length.

For the last assertion, assign each missing term to one of its
incident paid cubes. For one paid cube its possible roots lie
in a fixed collar, and the domain sum through any such root is
bounded by the strict entropy gap. Counting possible pointings
uses the spare regular rate just as above. A remainder has one
term per domain. The uniform bound per paid cube gives the cost
stated in the proposition.
\end{proof}

The cube at depth $s$ has side $M a_s$, where $M=L^{m_0}$ and
$a_s=L^{-s}$. On $\ell=2L^m$, $m\ge m_0$, the number of these cubes
in each direction is $2L^{m-m_0+s}$. This gives even counts on the
whole cofinal sequence. Every fixed board with side a power of
$L$ and integral in bare lattice units fits in all sufficiently
large members of the sequence. Domains at each of the actual
scales $0\le s\le K$ have the indexing used in the rooted counts.
The one-layer enlargement, including the fixed bare-barycentre
reach, stays within a fixed number of neighboring cubes since
$a\le a_s$. Thus the support-range and multiplicity conditions
in the subsequent geometric estimates refer to the actual energy
sum. If a filing uses additional domain types, their multiplicity
must be paid in the same entropy estimate.

The physical normalized integral at a stop is
$\widetilde Z(W)=e^{-V-\mathcal E(W)} Z(W)$. This uses the exact
history mass in \eqref{si:eq-frozen-transport}. On two admitted
common prefixes the same local formation pulls back the
same bare plaquette arrays. The equality of the corresponding
source energies and scalar shares on the common reads follows
from their definitions. Under an admitted change of lattice units
this statement uses the relabeling of the frozen tests, measures,
carriers and pointing data of \cite{RT}. At a stop, the natural
source norm is $\|w\|_{\flat,K-n}$. With unchanged physical units
the site count is $\ell^4 h_n^{-4}$, and with the units of the stop
it is the fourth power of the dilated side. These counts agree.
The upper and lower bounds on this normalized integral require
the retained-integral estimates and the unsourced floor.

If this same real run is admitted through a later level $K$,
write $\mathcal E_{(n)}(W)=\mathcal E(W)$ for the source sum
through its stop $n$. The two vacuum conventions are related by
the exact identity
\begin{equation}\label{si:eq-stopped-vacuum-conversion}
 \begin{split}
 e^{-E_{(K)}-\mathcal E_{(n)}(W)}Z(W)
 &=e^{E_{(n)}-E_{(K)}}
       e^{-E_{(n)}-\mathcal E_{(n)}(W)}Z(W).
 \end{split}
\end{equation}
Assume the full-level scalar convention and the bounded remaining
scalar prices of \cite[Proposition~2.5]{RT}, with
$\theta_j\ge1$ and $\theta_i-\theta_j\ge-C_{\rm exc}$ for $i\le j$.
That proposition bounds the positive multiplier by
\[
 e^{-C N_n}\le e^{E_{(n)}-E_{(K)}}
              \le e^{C(1+\log\theta_n)N_n},
 \qquad N_n=\ell^4h_n^{-4}.
\]
Thus both normalizations give an upper bound and a zero-source
floor with the corresponding adjusted stop-scale exponents.
The source sum still ends at $n$ in both conventions.


\subsection{The neighborhood where a partition logarithm exists}
\label{si:subsec-local-log}

An upper bound for a holomorphic partition function and a positive
value at zero do not exclude zeros throughout its source ball. The
following local statement is enough to define its derivatives at
zero when a logarithm is used.

\begin{samepage}
\begin{lemma}[A local logarithm]\label{si:lem-local-log}
Let $H$ be holomorphic on every finite-dimensional restriction of a
complex norm ball $\|W\|<R$, with compatible values on overlaps.
Suppose $|H(W)|\le M$ there and $H(0)\ge m>0$, where $M\ge m$.
On the original ball,
\[
 |H(W)/H(0)-1|\le (1+M/H(0))\|W\|/R.
\]
Moreover $H$ has no zero on
\[
           \|W\|<\frac{mR}{M+m}.
\]
The logarithm of $H(W)/H(0)$ defined to be zero at the origin is
holomorphic on every finite-dimensional restriction of this ball.
On the smaller ball $\|W\|<mR/(2M+m)$, its absolute value is at
most $\log2$.
\end{lemma}
\end{samepage}
\begin{proof}
Schwarz's lemma applied to $H-H(0)$ on each radial disc gives
the first estimate, since its absolute bound is $M+H(0)$.
For $W\ne0$, apply Cauchy's coefficient bound to the one-variable
function $\zeta\mapsto H(\zeta W/\|W\|)$ on $|\zeta|<R$.
Summing its Taylor series away from the constant term gives
\[
 |H(W)-H(0)|\le\frac{M\|W\|}{R-\|W\|}<m.
\]
Thus $|H(W)/H(0)-1|<1$, proving nonvanishing. The logarithm series converges locally
uniformly on each finite-dimensional restriction, and defines
compatible functions there. The value at zero is prescribed by
the same series. On the smaller ball the displayed difference
is less than $m/2$, so the ratio differs from one by less than
$1/2$. The absolute sum of its logarithm series is then at most
$\sum_{k\ge1}2^{-k}/k=\log2$.
\end{proof}

\begin{samepage}
\begin{corollary}[Finite-volume residual cumulants]
\label{si:cor-residual-cumulants}
Let $c_i>0$, and let $h$ be holomorphic on
$\sum_i c_i|z_i|<1$, with $h(0)=1$ and $|h(z)|\le Q$, where
$Q\ge1$. If $\mathcal E$ is holomorphic near zero,
$\Gamma(z)=e^{\mathcal E(z)}h(z)$ and $\mathcal E(0)=0$, then
\[
 \left|\partial_1\cdots\partial_n
       (\log\Gamma-\mathcal E)(0)\right|
   \le(\log2)\,[n(2Q+1)]^n\prod_i c_i.
\]
The logarithms are the branches equal to zero at the origin.
\end{corollary}
\end{samepage}
\begin{proof}
Apply \cref{si:lem-local-log} with $m=1$, $M=Q$ and $R=1$ in
the weighted parameter norm. In its smaller ball of radius
$1/(2Q+1)$, $\log h$ is
bounded by $\log2$. Choose coordinate circles with radii tending
upward to $1/[n(2Q+1)c_i]$. Their products remain inside that ball
before the limit. Multivariable Cauchy gives the product of the
inverse radii. Each variable is differentiated only once.
\end{proof}

For a normalized sourced partition function satisfying
$|H(W)|\le e^{C_+\ell^4}$ and $H(0)\ge e^{-C_-\ell^4}$, the
radius supplied by this lemma may decay exponentially in the
volume. No uniform global logarithm is asserted. The common cell
normalizers used below are instead built from the local source
energies themselves; their construction does not take a logarithm
of the full partition function.

\begin{samepage}
\begin{lemma}[Separated derivatives of local energies]
\label{si:lem-separated-energy}
Let $a_s=L^{-s}$ for $0\le s\le K$. Suppose a family of domains
$X$ has an integer size $m(X)\ge0$ and source dependence region
$X^+$ with
$\operatorname{diam}X^+\le C_0 a_s(m(X)+1)$. Let
$e_{s,X}(\sum_i z_i f_i)$ be holomorphic on
$\sum_i|z_i|\|f_i\|_\sharp<r$, with bound
$A e^{-\kappa m(X)}$, and depend only on the source in $X^+$.
Assume that the number of size-$m$ domains whose dependence
region meets $\operatorname{supp}f_1$ is at most
$H a_s^{-4} e^{\kappa_D m}$, where $\kappa>\kappa_D$.
Assume the domain sums converge normally on smaller source balls
at each fixed cutoff. For $n\ge2$ and pairwise support distance
at least $d>0$,
\begin{equation}\label{si:eq-separated-energy}
 \left|\left.\partial_1\cdots\partial_n
            \sum_{s=0}^K\sum_X e_{s,X}\!\left(\sum_i z_i f_i\right)
                  \right|_{z=0}\right|
 \le C A H\max\{1,d^{-4}\}(n/r)^n\prod_i\|f_i\|_\sharp.
\end{equation}
The constant depends only on $L,C_0$ and the entropy gap, and is
independent of $K$.
\end{lemma}
\end{samepage}
\begin{proof}
A zero test gives a zero derivative, so assume all test norms are
positive. Locality makes the mixed derivative vanish unless
$X^+$ meets every support. A contributing domain therefore has
$m(X)\ge d/(C_0a_s)-1$. Cauchy circles with radii approaching
$r/(n\|f_i\|_\sharp)$ bound its derivative by
$A(n/r)^n e^{-\kappa m(X)}\prod_i\|f_i\|_\sharp$.
The count through the first support and the strict entropy gap
bound the sum at scale $s$ by the same prefactor times
$C H a_s^{-4}e^{-c d/a_s}$.

For $0<d\le1$, the scales with $a_s\ge d$ form a geometric
sum bounded by $C d^{-4}$. At smaller scales write
$a_s\asymp d L^{-j}$; their sum is at most
$C d^{-4}\sum_{j\ge0}L^{4j}e^{-c' L^j}$, which converges.
For $d>1$, replacing $d$ by one only increases the sum.
Normal convergence justifies termwise differentiation at each
cutoff, and these estimates give its uniform bound.
\end{proof}

A finite-volume root count may take $H=C\ell^4$. A count anchored
to a fixed radius-$R$ support may instead take $H=C(R+1)^4$.
The same proof covers both uses. A point counterterm
$U(W)=\sum_p u_p(W(x(p)))$ has zero mixed derivative when the
supports are disjoint, since no point reads two of the variables.
Together with \cref{si:cor-residual-cumulants},
\cref{si:lem-separated-energy} gives a cutoff-independent separated
cumulant bound whenever $Q$ has such a bound at the chosen finite
volume. This remains a finite-volume statement if $Q$ depends on
$\ell$. A comparison of successive cutoffs is sufficient under the following
additional hypotheses.

\begin{samepage}
\begin{proposition}[A sufficient comparison of two cutoffs]
\label{si:prop-cutoff-comparison}
Fix the physical torus and real tests $f_1,\ldots,f_n$ whose
supports have mutual distance on the torus at least $d>0$, where $n\ge2$.
Use the same coefficients $z=(z_1,\ldots,z_n)$ for these tests
at every cutoff, on one fixed open polydisc about zero inside
all the admitted source balls. Write
\[
 \Gamma_K(z)=\frac{Z_K(\sum_i z_i f_i)}{Z_K(0)}
      =e^{\mathcal E_K(z)}h_K(z),\qquad h_K(0)=1.
\]
Assume $\mathcal E_K$ consists of the local source energies and
possibly point terms, as in \cref{si:lem-separated-energy}.
Use uniform locality, decay and rooted counts on that polydisc
at the physical scales $a_s=L^{-s}$, including a strict entropy
margin. Suppose the two energies at cutoffs $K$ and $K+1$ can
be paired for all indices at $s\le K/2$, padding absent terms
by zero, using domains with the same complete source dependence
regions and the same rooted count. Assume
that for each paired index the difference is bounded by
\[
 C(1+s)\epsilon_K e^{-\kappa m(X)},
\]
with the same strict entropy margin. On the common polydisc
suppose $|h_K-1|,|h_{K+1}-1|\le1/2$ and
$|h_{K+1}-h_K|\le\epsilon_K$. If $\sum_K\epsilon_K<\infty$,
the separated cumulants for this tuple converge along the
whole cutoff sequence.
\end{proposition}
\end{samepage}
\begin{proof}
Point terms have zero mixed derivative because each sees at
most one of the separated tests. For a paired domain the
mixed derivative of its difference vanishes unless its
complete source read meets every support. Its size therefore
satisfies $m(X)\ge d/(C_0a_s)-1$. Cauchy's estimate on a fixed
smaller polydisc and the same rooted count as in
\cref{si:lem-separated-energy} bound the derivative of the
coarse paired sum by
\[
 C(f_1,\ldots,f_n)\epsilon_K
       \sum_{s\ge0}(1+s)L^{4s}e^{-cdL^s}.
\]
At $s>K/2$ bound the two local energies separately by the same
individual estimates. The tail $\sum_{s>K/2} L^{4s}e^{-cdL^s}$
is summable in $K$, since interchanging the two sums gives at
most $\sum_s(2s+1)L^{4s}e^{-cdL^s}$.

For the residuals, the segment between $h_K(z)$ and $h_{K+1}(z)$
stays in the closed disc of radius $1/2$ about one. Integrate
$1/\zeta$ along it on the branch based at one to obtain
$|\log h_{K+1}-\log h_K|\le2\epsilon_K$. Cauchy's formula bounds
the mixed derivative by a constant times $\epsilon_K$. The
cumulant is the derivative of $\mathcal E_K+\log h_K$ at zero,
so its consecutive differences are summable. The triangle
inequality now gives convergence. The hypotheses concern paired
energies of the actual varying source and the common residual
disc. Individual bounds at the two cutoffs supply neither
comparison.
\end{proof}



\section{An upper bound for the sourced history sum}
\label{si:sec-history-stability}

Fix the cutoff $K$, torus side $\ell$, and a real admitted run.
The physical spacing at level $j$ is $L^{j-K}$ and its number of
sites is $N_j=\ell^4L^{4(K-j)}$. At a stop $n\le K$ let $E_{(n)}$
be its common vacuum number, with the sign convention of
\cite[Equation~(2.iv)]{RT}.
Write $Z(W)$ for the unnormalized bare integral with insertion
$e^{\mathcal O_a(W)}$, where $a=L^{-K}$ and $x(p)$ is the centre of
plaquette $p$. Every use of a source below is within the common
ordinary analytic radius of the frozen construction. On the tori
used for state transfer we take this radius to be $2r>0$ in the
physical norm of the next section.

Write $\mathcal E_{(n)}(W)$ for the common source number at this
stop. It consists of the local energies of
Proposition~\ref{si:prop-local-energies}, and any separately filed
point counterterm $U(W)$, holomorphic with $U(0)=0$. In that convention point terms are
counted outside the local energies. The native filing has $U=0$.

The quantitative inputs for the sourced sum concern the same
positive reference laws and the same record spaces as the real run.
The complete sourced densities can take complex values. We use the
following local conditions on these densities.

\begin{enumerate}
\item The exact frozen pushforward gives complex label measures
  $\tau_{\mathcal S}(W)$ with total $Z(W)$ after summing their
  masses. The product coordinates and their Borel domains are the
  actual ones. Their source-independent reference factors and the
  disjointness or normalized counting of hidden records are as in
  the zero-source history sum. Write $e^{v_{\mathcal S}(W)}$ for the
  numerical factor on label $\mathcal S$. Use the actual scalar
  assignment of \eqref{si:eq-vacuum-share}, with the missing shares
  charged at their original levels. If point terms are separately
  filed, their missing shares must have the same per-address bound.
  The charged addresses are in the live record or owned by a
  retained operation which has remade them. As in the loaded
  history estimate of \cite[Proposition~8.5]{RT}, the weight
  $v(\omega)$ of every record compatible with the visible label
  covers the own-scale sites assigned to it, with bounded
  multiplicity.
\item For the source-dependent changes of the remaining regular,
  remainder and boundary exponent, take the sum of absolute
  bounds on the actual carriers before summing the histories.
  On a final-site read the rooted bound is at most
  $A_0(1+\log x_n)$, including the contributions of all earlier
  levels after conversion to final-site units. For example, a
  sufficient set of separate origin bounds at age $h$ is
  $a_h L^{-4h}$ per origin site with
  $\sum_{h\ge0}a_h\le C(1+\log x_n)$. On the other reads the rooted bound is
  $A_1(1+\log x_j)$ per charged site at its original level $j$.
  These bounds apply on the same admitted real stored variables
  on which the sourced terms are evaluated. A read owned by a
  retained operation is charged to its stored sites. This input
  includes the terms attached to an island boundary and the
  factors carried outside a drop-out.
\item The actual retained operations satisfy the complete-read
  and fresh amplitude bounds of Lemma~\ref{si:lem-retained-source-excess},
  including the collar constants and the kept-response corrections.
  The resulting source-exponent changes obey
  $|\Psi_\omega(W)|\le A_2 n_\omega$ on the allowed exterior and
  stored reads. The real weighted masses with their exterior
  suprema inside the integrals are available at the fixed
  increased reserve weights used below, chosen before
  the large inverse-coupling floor. Each allocated failure factor
  on its real support has an absolute majorant $e^{-b}$ times a
  declared residual reference factor, with a positive uniform
  constant in the paying gain $b$. The needed real retained mass
  bounds include those residual factors and cuts. At the increased
  reserve they retain the claimed order of every mark gain spent
  in that mass. Otherwise an assigned last-test share must supply
  that gain separately. The remaining
  Wilson factor has nonnegative real plaquette coefficients after
  the share allocation. Coefficients read on different plaquette
  fields use their own allocations. The change $\Psi_\omega$
  concerns the remaining factors after the allocated shares and
  the separately charged local changes have been removed.
\end{enumerate}

These are estimates on local formation, allocation and retained
integrals. No sum over the final supports is among them.
For the ordinary old-origin regular part of the second input,
\eqref{si:eq-member-sum} gives a bounded cost per final window
when its actual reads satisfy that proposition's hypotheses.
Covering the final sites by a bounded number of such windows per
site gives the corresponding part of the final-site cost.
Boundary and retained reads still use their stated local bounds.
For an
action weight at one fixed plaquette field, ordered counterterm
profiles give the convex combination in the real history argument
of \cite[Equation~(2.vi)]{RT}. The same linear identity applies to the real parts
of complex weights. For example, if all running weights on that
plaquette differ from their positive real reference by at most a
quarter of that reference, its real action coefficient is at least
three quarters of the coefficient obtained from the reference
weights. A positive allocation has to be checked on the same
carrier before spending its gain.

\begin{samepage}
\begin{proposition}[Ordinary complex-source stability]
\label{si:prop-ordinary-stability}
Use the actual frozen source representation and the three local
inputs above. Assume the unsourced history geometry, real cost,
record and paying-event estimates of the upper history bound in
\cite[Proposition~8.4]{RT}, with actual rounded running-minimum
radii of exponent $r_{\mathrm{cell}}\ge1$ and paying gains of
order $(\log\theta_j)^p$ for $p>6r_{\mathrm{cell}}$. The frame
has $\theta_j\ge\theta_n-c_*$ and $x_j\asymp\theta_j$ for $j\le n$.
Then, above a fixed inverse-coupling floor,
\begin{equation}\label{si:eq-ordinary-upper}
 |Z(W)|\le
 \exp\{E_{(n)}+\operatorname{Re}\mathcal E_{(n)}(W)
                  +C(1+\log x_n)N_n\}.
\end{equation}
Under the actual empty-label assumptions of the common-normalized
partition bounds in \cite[Theorem~E]{RT}, also
\begin{equation}\label{si:eq-unsourced-floor}
              Z(0)\ge \exp\{E_{(n)}-C(1+\log x_n)N_n\}.
\end{equation}
The constants are uniform in cutoff, volume, admitted frame and
source, with the displayed data fixed. At a terminal first-passage
datum $n=K$ with $x_K$ in the fixed window about $g^{-2}$, the
error terms in both estimates are at most $C_g\ell^4$.
\end{proposition}
\end{samepage}
\begin{proof}
Factor out $e^{E_{(n)}+\mathcal E_{(n)}(W)}$ from the sum of the
complex label masses. At zero the difference between the label
number and $E_{(n)}$ is the real vacuum price of the history bound.
For the sourced part, subtract the empty-label shares from the
shares in the label at hand. At each missing or changed root the
absolute source price is at most the convergent per-root sum of
Proposition~\ref{si:prop-local-energies}. A change of a share can
cost twice that bound. With the local multiplicities this adds at
most a constant per charged site at its original level. The second
input adds at most $A_0(1+\log x_n)N_n$ on final-site reads and changes the
own-scale charge constant on other reads.

For example, a spent energy share $Q$ can pay a gain $b$ if the
failure forces $Q\ge4b$ and the source perturbation of that share
satisfies $\operatorname{Re}\Delta\le Q/4$. On that support,
\[
  1_{\{\mathrm{failure}\}}\bigl|e^{-Q+\Delta}\bigr|
       \le 1_{\{\mathrm{failure}\}}e^{-3Q/4}
       \le e^{-b}e^{-Q/2}.
\]
The residual reference factor is then $e^{-Q/2}$ with the cut.
Any reference mass cost of retaining the reduced energy belongs
to its local mass bound. The source changes and the reassigned
prices of a retained operation are bounded by $c v(\omega)$ for
a fixed $c$, since $1+\log x_j\ge1$ in the weights on stored sites.
Increase the reserve by that $c$ while keeping the same restrictions
on its visible records. This comparison takes place on the stored
variables, before the exterior supremum and integration. Apply
the product majorant of \cite[Lemma~8.1]{RT} to the disjoint integration groups.
It allows coupling of their exterior arguments without exchanging
their suprema with their integrals. Use the same normalized
hidden-record reference or the same disjoint record sum as in the
real estimate. A hidden failure already used in a retained mass
contributes no second gain.

For a final component $C$, let $v_C$ be its charged volume and
$G_C$ the sum of its surviving paying gains. The resulting bound
has a common factor $e^{C(1+\log x_n)N_n}$ and a factor
$e^{A v_C-G_C}$ per component record. The own-scale genealogy
charge of \cite[Lemma~8.2]{RT} bounds $v_C$ by a fixed constant times
$\sum_e M^4\check R_{j(e)}^6$, summed over these same events.
The volume $v_C$ already counts $1+\log x_j$ per own-scale site.
The running-minimum comparison gives $1+\log x_j\le C\check R_j$
and $\check R_j\le C(1+\log\theta_j)^{r_{\mathrm{cell}}}$.
Since $p>6r_{\mathrm{cell}}$, increasing the floor makes half of $G_C$ pay this
charge with any fixed additional entropy constant. The other
half leaves a gain at least $c(1+\log\theta_n)^p$, because every
nonempty component has a paying event and its level $j$ satisfies
$\theta_j\ge\theta_n-c_*$. In particular, for a fixed sufficiently large
$\lambda$ and the rooted code constant $c_{\mathrm e}$ of the
real estimate,
\[
   e^{A v_C-G_C}\le
       e^{-c(1+\log\theta_n)^p-\lambda|C|}e^{-c_{\mathrm e}v_C}.
\]
The sum over records of $e^{-c_{\mathrm e}v_C}$ is at most one.
Discard the remaining positive event gain when summing the
partition function. It has not been converted to a probability.
On a final graph of degree at most $D_0$, at most
$D_0^{2(m-1)}$ connected sets of size $m$ are rooted at a given
cell, by traversing a spanning tree twice. Thus the sum of
$e^{-\lambda|C|}$ over all possible final components is at most
$C N_n$ for $\lambda>2\log\max(2,D_0)$. Bounding the collection sum
by the product of $1+e^{-\lambda|C|}$ over these components gives
at most $e^{C N_n}$. This proves \eqref{si:eq-ordinary-upper}.
Apply the zero-source lower bound of \cite[Theorem~E]{RT} for
\eqref{si:eq-unsourced-floor}. Here $\mathcal E_{(n)}(0)=0$, so its
vacuum factor is exactly $E_{(n)}$. At $n=K$, $N_K=\ell^4$, and
the first-passage window bounds $1+\log x_K$ at fixed $g$.
\end{proof}

\begin{samepage}
\begin{corollary}[Finite-torus separated cumulants]
\label{si:cor-finite-torus-cumulants}
On a compatible torus of side $\ell\ge1$, assume the ordinary
source upper bound and zero-source floor with errors
$C_+\ell^4$ and $C_-\ell^4$, respectively, uniformly on the
source ball of radius $2r$. Assume the domain terms of their
common source number $\mathcal E$ obey the local-energy estimates
of \cref{si:prop-local-energies}, with the actual cubical indexing
and strict entropy margin of \cref{si:lem-source-count}.
A separate holomorphic point counterterm vanishing at zero is
allowed in this same source number. For $q\ge2$, let
$f_1,\ldots,f_q$ be real Lipschitz
torus tests whose supports have mutual torus distance at least
$d>0$. Their joint cumulant under the finite Wilson probability
satisfies
\[
 |\kappa_q^{\mu}(\mathcal O_a(f_1),\ldots,\mathcal O_a(f_q))|
 \le C_q(d)e^{q c_V\ell^4}\prod_i\|f_i\|_\sharp .
\]
The constants are uniform in cutoff and support positions when
the input constants are uniform. No support-radius bound is
required, and $c_V$ is independent of $q$.
\end{corollary}
\end{samepage}
\begin{proof}
If any test is zero, the cumulant is zero. Assume all their
norms are positive. Write $\Gamma=Z(W)/Z(0)$ and
$h=e^{-\mathcal E(W)}\Gamma$.
The two partition bounds give $h(0)=1$ and
$|h|\le Q=e^{(C_++C_-)\ell^4}$. Increase the nonnegative cost
constants if necessary. Apply \cref{si:cor-residual-cumulants}
with $c_i=\|f_i\|_\sharp/r$. It bounds the mixed derivative of
$\log h$ by $C_qQ^q\prod_i\|f_i\|_\sharp$.
For the derivative of $\mathcal E$, use
\cref{si:lem-separated-energy} with the global root count
$H=C\ell^4$. Its contribution is at most
$C_q\ell^4\max(1,d^{-4})\prod_i\|f_i\|_\sharp$.
Point terms contribute zero to this mixed derivative by separation.
The cumulant is the sum of these two derivatives. The inequality
$\ell^4\le e^{\ell^4}$ absorbs the polynomial volume factor;
one may take $c_V=C_++C_-+1$ after adjusting $C_q(d)$.
\end{proof}

\section{Transfer at fixed spacing}\label{si:sec-state-transfer}

Fix a cutoff--coupling datum $a=L^{-K},\beta$ and a cofinal admitted
family of tori of side $\ell=2L^m$. Hold $a,\beta$ fixed in the
volume limit. Assume ordinary source bounds of the form
\eqref{si:eq-ordinary-upper} and \eqref{si:eq-unsourced-floor}, with
errors $C_+\ell^4$ and $C_-\ell^4$. The constants can be uniform
over a class of data satisfying the same source estimates. The
preceding terminal first-passage construction is one sufficient
case when its local hypotheses hold.

Write $V_\ell$ for the real vacuum number in these bounds, and
write their common source number as
\[
 \mathcal E_\ell(W)=U(W)+\sum_{s=0}^K\sum_D e_{s,D}(W),
 \qquad a_s=L^{-s}.
\]
A separately filed point term $U(W)=\sum_p u_p(W(x(p)))$ uses
holomorphic functions vanishing at zero and locally bounded
uniformly in $p,\ell$ at each fixed cutoff. Each domain term
vanishes at zero, is holomorphic and depends on $W$ only through
its values on $D^+$. Assume
$|e_{s,D}|\le A e^{-\kappa m(D)}$ and
$\operatorname{diam}D^+\le C_0 a_s(1+m(D))$, with the same
nonnegative integer size $m(D)$, and assume normal convergence on
smaller source balls. The count needed for ordinary transfer is
\begin{equation}\label{si:eq-ordinary-root-tail}
 \sum_{\substack{D^+\cap B_R\ne\varnothing\\m(D)\ge t}}
 e^{-\kappa m(D)}
 \le C(R+1)^4a_s^{-4}e^{-ct}
 \qquad(t\ge0)
\end{equation}
uniformly in every radius-$R$ ball $B_R$ and the torus. The
rooted domain count of \cref{si:lem-source-count} implies
\eqref{si:eq-ordinary-root-tail} because $a_s\le1$. Its finer
dependence on $R/a_s$ and its global estimate will be used in
the graded argument.

Use the physical test norm $\|\cdot\|_\sharp$ of
\eqref{si:eq-source-norm} at $s=0$. The ordinary source bound is used for all finite
representations $W=\sum_i z_i f_i$ on the torus with
$\sum_i|z_i|\|f_i\|_\sharp<2r$. Constant sources can be included
among these real profiles. Every local term depends on the function
restricted to its read region. In particular its value does not
change under a different admitted list representing that same
restriction. The transfer uses \eqref{si:eq-ordinary-upper} and
\eqref{si:eq-unsourced-floor} with constants $C_+$ and $C_-$ per
unit volume, respectively.

Fix a board of side $l=L^b\ge a$ for an integer $b$, with faces
bisecting links, as in \cite[Theorem~C]{LP}. Let
$p_{l,\ell}(F)$ be its finite chessboard seminorm. It takes the
$N^{-1}$ power of the expectation of the disseminated product,
where $N=(\ell/l)^4$ and all copies of odd reflection parity are
conjugated. The expectation in this definition is for the original
positive Wilson measure at zero source. Suppose real test functions
$f_i$ assigned to a particular cell have supports in $m_C$ balls
of radius $R$, at distance at least $d_{\mathrm f}>a$ from its
faces. Put $w_C=\sum_i z_if_i$. The plaquette terms in
$\mathcal O_a(w_C)$ use only bonds in the open cell. In every
dissemination, the even and odd copy profiles of each $f_i$ are
real and vanish near the faces. Their Lipschitz constants on the
whole torus are at most $\operatorname{Lip}f_i$. Indeed the copies
have that bound on each cell, agree at the faces, and the bound on
a geodesic follows by adding its cell segments. Their suprema
are also bounded by $\|f_i\|_\infty$. Hence the disseminated source
$W^{\mathrm{ch}}$ has an ordinary representation of radius at
most $2\sum_i|z_i|\|f_i\|_\sharp$. This represents it with the two
lists $z_i$ and $\bar z_i$ evaluated after the ordinary analytic
functions have been defined.

\begin{samepage}
\begin{lemma}[A holomorphic cell normalization]
\label{si:lem-cell-normalizer}
With these hypotheses and $\sum_i|z_i|\|f_i\|_\sharp<r$, there
is a holomorphic function $E_{C,\ell}(z)$ with $E_{C,\ell}(0)=0$
such that
\begin{equation}\label{si:eq-cell-normalizer}
 p_{l,\ell}(e^{\mathcal O_a(w_C)})
     \le \exp\{\operatorname{Re}E_{C,\ell}(z)+B_C\},
 \qquad
 B_C=(C_++C_-)l^4+C m_C(R+1)^4
       \sum_{s\ge0}a_s^{-4}e^{-c d_{\mathrm f}/a_s}.
\end{equation}
One may replace the second term in $B_C$ by any uniform bound
on the absolute discarded source numbers per board cell.
At fixed cutoff the functions $E_{C,\ell}$ are locally bounded
uniformly in $\ell$. Along a cofinal subsequence they converge
locally uniformly to a holomorphic $E_C$. For finitely many chosen
cell profiles one subsequence suffices and depends only on the
finite-torus source numbers and the profiles.
\end{lemma}
\end{samepage}
\begin{proof}
The disseminated product is $e^{\mathcal O_a(W^{\mathrm{ch}})}$,
since the plaquette action is invariant under orientation reversal.
Its expectation is the nonnegative real ratio
$Z(W^{\mathrm{ch}})/Z(0)$, by finite chessboard positivity.
The source bounds therefore control the $N$th root of its modulus.

In the common source number keep the point terms and the domain
terms with their complete read regions inside one open cell.
Each kept term can be evaluated on the source copy in that cell
alone. The kept terms in even cells are holomorphic in $z$.
An odd cell term has the form $g(\bar z)$ for a holomorphic
$g$. Use its conjugate $\overline{g(\bar z)}$ when forming the
average. This is holomorphic by its local power series, and its
real part is that of the actual term. Divide the resulting sum by
$N$ to define $E_{C,\ell}$. The sum vanishes at zero.

A discarded term is zero if it reads no nonzero source value.
Otherwise its read region meets a support ball and is not contained
in that ball's cell. Its diameter is at least $d_{\mathrm f}$,
even on the periodic quotient. Thus $m(D)\ge d_{\mathrm f}/(C_0a_s)-1$.
Use \eqref{si:eq-ordinary-root-tail} once for each of the $N m_C$
ball copies. The exponential estimate for $e_{s,D}$ yields the
second term of \eqref{si:eq-cell-normalizer} after division by
$N$. No nonzero point term is discarded. The series is finite
because $a_s^{-1}=L^s$. The extensive part of the partition bound
divided by $N$ is $(C_++C_-)l^4$, which proves the estimate.

At fixed cutoff the point terms per cell are bounded on smaller
source balls by their fixed-cutoff bounds times $C l^4a^{-4}$.
For the kept domain terms use \eqref{si:eq-ordinary-root-tail} with
$t=0$ and the finite sum over $s\le K$. This gives local bounds
independent of $\ell$ on each smaller source ball. Montel's theorem
gives a cofinal subsequence, and a successive subsequence choice
works for finitely many cells. The limit has the same zero value.
\end{proof}

\begin{samepage}
\begin{theorem}[One source normalization for all periodic states]
\label{si:thm-state-transfer}
Assume the ordinary upper and zero-source lower bounds on the
cofinal tori, the local source numbers, diameter bound and tail estimate
\eqref{si:eq-ordinary-root-tail} above,
and fixed-spacing chessboard domination of \cite[Theorem~C]{LP}.
Let $f_1,\ldots,f_q$ be real Lipschitz tests with supports in
radius-$R$ balls in physical space. For this list and datum there
is a holomorphic function $E(z)$ on
$\sum_i|z_i|\|f_i\|_\sharp<r$, with $E(0)=0$, such that for every
periodic Wilson DLR state $\rho$ at the datum,
\begin{equation}\label{si:eq-state-generating}
 \left|e^{-E(z)}\rho\left(e^{\sum_i z_i\mathcal O_a(f_i)}\right)\right|
 \le Q_{q,R}.
\end{equation}
The same $E$ serves all such states. The constant $Q_{q,R}$ is
uniform in cutoff, states and support positions. For one test
$f$ of radius $R$ and any two states $\rho,\nu$,
\begin{equation}\label{si:eq-state-mean-difference}
       |\rho(\mathcal O_a(f))-\nu(\mathcal O_a(f))|
            \le C(R+1)^4\|f\|_\sharp .
\end{equation}
\end{theorem}
\end{samepage}
\begin{proof}
Omit zero tests. Fix $d_{\mathrm f}>1$. In each coordinate the
forbidden board offsets for one support ball are contained in an
interval of length $2(R+d_{\mathrm f})$ modulo the side, with a
rounding cost of at most $2a$ on the link-bisecting grid. Choose the
allowed side $l$ to be the first power of $L$ exceeding
$4q(R+d_{\mathrm f}+1)$. The forbidden offsets cannot cover the
available grid in any coordinate. This puts every ball at distance
at least $d_{\mathrm f}$ from its cell faces and uses
$l\le 4Lq(R+d_{\mathrm f}+1)$. The torus sequence eventually admits
that board with even counts.

Group the tests by their cells and apply
Lemma~\ref{si:lem-cell-normalizer} to each occupied cell. Use one
common subsequence and put $E=\sum_C E_C$. At fixed cutoff the
cell observables are bounded functions of finitely many bonds,
and their exponentials are bounded complex cell functions at
every fixed $z$. Chessboard domination holds for every periodic
DLR state along every compatible cofinal sequence. Applying it
along this one gives \eqref{si:eq-state-generating}, with
$Q_{q,R}=\exp(\sum_C B_C)$ bounded using at most $q$ cells and the
side bound. The subsequence was chosen before using any state.
This argument does not select the state as a torus weak limit.

For one test put
\[
 h_\rho(z)=e^{-E(z)}\rho(e^{z\mathcal O_a(f)}),\qquad h_\rho(0)=1.
\]
Holomorphy of the expectation follows from boundedness
of the observable at fixed spacing. Its derivative is
$h_\rho'(0)=\rho(\mathcal O_a(f))-E'(0)$.
Cauchy's estimate gives $|h_\rho'(0)|\le Q_{1,R}\|f\|_\sharp/r$.
Subtract the corresponding identity for $\nu$. To obtain the
stated radius dependence, use a smooth partition of unity by
integer translates, supported in fixed-radius balls with uniformly
bounded first derivatives. Such a partition is obtained by dividing
a nonnegative compactly supported bump, positive on a unit cube,
by the positive periodic sum of its translates. At most
$C(R+1)^4$ pieces meet the support of $f$. Each piece times $f$
has $\sharp$ norm at most $C\|f\|_\sharp$ and fixed support radius,
so summing the fixed-radius mean bounds proves
\eqref{si:eq-state-mean-difference}. The derivative of the common
normalizer itself can grow with the cutoff.
\end{proof}

\begin{samepage}
\begin{corollary}[Separated cumulants]\label{si:cor-state-cumulants}
Under Theorem~\ref{si:thm-state-transfer}, let $q\ge2$.
If the supports have pairwise distance at least $d>0$, then
for every periodic state at the datum,
\[
 \bigl|\kappa_q^\rho(\mathcal O_a(f_1),\ldots,\mathcal O_a(f_q))\bigr|
 \le C_q(d,R)\prod_{i=1}^q\|f_i\|_\sharp.
\]
For arbitrary supports, the difference of these cumulants in two
states has a bound with a constant depending on $q,R$ alone.
The constants are cutoff independent.
For any two such states $\rho,\nu$ and pairwise separated tests,
the mixed moment centered at $\nu$ obeys
\[
 \left|\rho\prod_{i=1}^q
       (\mathcal O_a(f_i)-\nu(\mathcal O_a(f_i)))\right|
 \le \left(\sum_{\pi}\prod_{I\in\pi} b_{|I|}\right)
                         \prod_i\|f_i\|_\sharp,
\]
where $b_1=C(R+1)^4$, $b_m=C_m(d,R)$ for $m\ge2$, and the sum
runs over the set partitions of the indices.
\end{corollary}
\end{samepage}
\begin{proof}
Write $h_\rho(z)$ for the holomorphic function inside the modulus
in \eqref{si:eq-state-generating}. Apply \cref{si:cor-residual-cumulants} to $h_\rho$,
with $c_i=\|f_i\|_\sharp/r$ and $Q=Q_{q,R}$. Zero tests can be
omitted. It bounds the mixed derivative of $\log h_\rho$ by
\[
 (\log2)\,[q(2Q_{q,R}+1)/r]^q\prod_i\|f_i\|_\sharp.
\]
The common derivative of $E$ cancels in a difference of cumulants
of two states. Applying this bound to each state gives the
overlapping-support assertion with a cutoff-independent constant.

For one state, it remains to bound the mixed derivative of $E$.
If several cells are occupied, each cell summand lacks at least
one variable. If only one cell is occupied, point terms have zero
mixed derivative by separation. A contributing retained domain
must see every support, so its diameter is at least $d$. Its size is at least $d/(C_0 a_s)-1$. Cauchy's formula with
one circle for each of the $q$ source variables gives the
per-domain bound
$A(q/r)^q e^{-\kappa m(D)}\prod_i\|f_i\|_\sharp$.
Apply \eqref{si:eq-ordinary-root-tail} through the first support
ball in each copy, and divide by the number of copies. The scale
sum is $\sum_s a_s^{-4}e^{-c'd/a_s}\le C\max(1,d^{-4})$,
by the split at $a_s=d$ used in the proof of
\cref{si:lem-separated-energy}. Thus the mixed derivative of
$E_{C,\ell}$ is bounded by
\[
 C(R+1)^4\max(1,d^{-4})(q/r)^q\prod_i\|f_i\|_\sharp .
\]
Local uniform convergence passes it to $E_C$.
Adding the log bound proves the cumulant estimate. The moment
formula sums products of cumulants over set partitions. In a
singleton block use \eqref{si:eq-state-mean-difference}. All other
cumulants are unchanged by the constant shifts, so the preceding
bound applies to each of those blocks.
\end{proof}

\begin{samepage}
\begin{corollary}[Local differences between states]
\label{si:cor-local-state-differences}
Use Theorem~\ref{si:thm-state-transfer} and the lattice translation
invariance of the periodic states. Let $a\le d\le d_0$, with
$d_0>0$ fixed and sufficiently small, and let the supports of
$f_1,\ldots,f_q$ lie in one ball of radius $5d$. Write $D_q$ for
the difference of their order-$q$ cumulants in two periodic
states at the datum, viewed as a multilinear form of the tests.
For $q=1$ use the difference of the means. Then
\[
 |D_q(f_1,\ldots,f_q)|\le C_q d^4\prod_i\|f_i\|_\sharp.
\]
If $\operatorname{Lip} f_i\le\|f_i\|_\infty/(\vartheta d)$ for
a fixed $\vartheta>0$, the right side is at most
$C_{q,\vartheta}d^{4-q}\prod_i\|f_i\|_\infty$.
\end{corollary}
\end{samepage}
\begin{proof}
Use lattice shifts in a grid of spacing $a\lceil12d/a\rceil$,
restricted to a fixed cube of side $1/2$. Their number $P$ is
at least $c d^{-4}$. The translated balls are disjoint and fit in
a single unit ball about the original centre for small $d_0$.
For each argument $i<q$ take independent signs $\epsilon_v^i$ on
these shifts, and put $\epsilon_v^q=\prod_{i<q}\epsilon_v^i$.
Write $F_i=\sum_v\epsilon_v^i f_i(\cdot-v)$. At most one copy is
nonzero at any point. For points in different copy balls, split
their connecting segment at the two ball boundaries, where the
copies vanish. The lengths of the two inner segments sum to at
most the full distance. This gives $\|F_i\|_\sharp\le\|f_i\|_\sharp$.

In the multilinear expansion, averaging the signs retains precisely
the terms with all shifts equal, since the average is
$\prod_{i<q}\mathbf1_{v_i=v_q}$. Simultaneous lattice translation
leaves both states' cumulants invariant. Thus the average is
$P D_q(f_1,\ldots,f_q)$. For $q=1$ the same formula holds with all
signs one. Apply the overlapping-support bound of the preceding
corollary in the unit ball, or \eqref{si:eq-state-mean-difference}
when $q=1$, and divide by $P$. The final assertion follows from
the definition of $\|\cdot\|_\sharp$ and $d\le1$.
\end{proof}


\section{Separated moments with a fixed factorial exponent}
\label{si:sec-graded}

\begin{samepage}
\begin{lemma}[One-cell cost across scales]\label{si:lem-scale-cell-price}
Assume the ordinary partition bounds, source locality and tail
estimate of Lemma~\ref{si:lem-cell-normalizer}. Assume that this
source system has the actual cubical domain indexing and complete
fixed-layer reads of Lemma~\ref{si:lem-source-count}. Require
$|e_{s,D}(W)|\le A e^{-\kappa m_s(D)}$ for its actual tree length,
with $\kappa>\kappa_D$. Let a board have side
$l\le1$. Let a real test $v$ be supported in a ball of radius
$\delta$ inside one cell, at distance at least $\delta$ from the
faces, with $16a\le\delta\le l/4$. There are constants $Q_0,p$
independent of the cutoff, cell and torus such that
\begin{equation}\label{si:eq-scale-cell-price}
 p_{l,\ell}\!\left(e^{z\mathcal O_a(v)-E_{v,\ell}(z)}\right)
       \le Q_0(l/\delta)^p,\qquad |z|\|v\|_\sharp<r.
\end{equation}
The normalization is the one in Lemma~\ref{si:lem-cell-normalizer}.
It has the same fixed-cutoff locally uniform subsequential limits.
\end{lemma}
\end{samepage}
\begin{proof}
All contributing plaquettes use bonds of the open cell, by the
distance to its faces. The reflected source is admitted in the
ordinary radius $2r$, using the real reflected profiles in even
and odd cells and the coefficients $z$ and $\bar z$, respectively.
Each profile vanishes near the cell faces. The common source number
is split into kept and discarded terms as in
Lemma~\ref{si:lem-cell-normalizer}. A nonzero discarded term sees a
source ball and its complete read is not contained in that ball's
cell. Thus $m_s(X)\ge c_0\delta/a_s-C_0$.

At scales $a_s\ge l$, bound all terms by the global sum in
\eqref{si:eq-source-count} and divide by $(\ell/l)^4$. This costs
$C(l/a_s)^4$ per level. At $\delta\le a_s<l$, sum the ball count
over the source copies and divide by their number. The cost per
level is at most $C$, and there are at most $1+\log_L(l/\delta)$
levels. At $a_s<\delta$ the rooted tail gives cost at most
$C(1+\delta/a_s)^4e^{-c\delta/a_s}$. The coarse and fine costs have
uniformly bounded sums. For the fine sum, $\delta/a_s$ grows by
$L$ each step and the polynomial can be absorbed in half the
exponential. The partition cost is $Cl^4\le C$, and every
nonzero point term is kept. Taking the exponential of the
resulting bound $C(1+\log(l/\delta))$ gives the assertion.
\end{proof}

\begin{samepage}
\begin{theorem}[Graded separated moments]\label{si:thm-graded-moments}
Assume the ordinary source estimates and actual rooted domain
indexing of Lemma~\ref{si:lem-scale-cell-price}. Let
$f_1,\ldots,f_n$ be Lipschitz tests, each supported in a ball of
radius $R$. For $n\ge2$ assume that their supports have mutual
distance at least $d>0$. Set $\bar d=\min(d,1)$, with $\bar d=1$
for $n=1$. Let $\rho,\nu$ be any two periodic Wilson Gibbs states
at the same bare datum. Then there are a fixed integer $k$ and
$B(R)=C(R+1)^4$ such that
\begin{equation}\label{si:eq-graded-moment}
 \left|\rho\!\left(\prod_{i=1}^n
       [\mathcal O_a(f_i)-\nu(\mathcal O_a(f_i))]\right)\right|
    \le B(R)^n\bar d^{-5n}(n!)^k\prod_{i=1}^n\|f_i\|_\sharp.
\end{equation}
For $n\ge2$ the joint cumulant of the observables under $\rho$
has the same bound with $2B(R)$ in place of $B(R)$. The estimates
also hold if both expectations are the Wilson expectation on a
compatible admitted finite torus, with distances measured on that
torus. The constants are uniform in $n$, cutoff, volume, state
and support positions. They may depend on the fixed source
estimates and geometry.
\end{theorem}
\end{samepage}

\begin{proof}
We first take real tests. Choose a power of $L$ for the cell side,
$\bar d/(16L)<l\le\bar d/16$, and put $\delta=l/(64n)$.
If $a>\delta/16$, counting plaquette centres gives
\[
 \|\mathcal O_a(f_i)-\nu(\mathcal O_a(f_i))\|_\infty
 \le C(R+1)^4a^{-4}\|f_i\|_\sharp
 \le C(R+1)^4\delta^{-4}\|f_i\|_\sharp.
\]
The product of these bounds is within the estimate proved below.

Now let $a\le\delta/16$. Use a smooth partition of unity on a
grid of step $\delta/8$, and write $f_i=\sum_\gamma v_{i,\gamma}$.
One construction uses nonnegative translates of a fixed bump and
divides by their sum, which is uniformly bounded below. The bumps
have bounded overlap, support in radius-$\delta$ balls and
Lipschitz constant $C/\delta$ after normalization. The grid
tiles a compatible torus since $\ell/l$ and $l/(\delta/8)$ are
integers. There are at most $C(1+R/\delta)^4$ nonzero pieces per
test, and $\|v_{i,\gamma}\|_\sharp\le
C\delta^{-1}\|f_i\|_\sharp$.

For any choice $v_i$ of one piece per test, choose a common board
offset. In each coordinate there are $l/a$ offsets of link faces.
For each ball at most $4\delta/a+2$ of them put some face within
distance $\delta$ of the ball. Since $n(4\delta/a+2)<l/a$, an
offset remains in each coordinate. The selected balls are inside
their cells with the required reserve. Nonzero pieces from
different tests lie in different cells, since their nonzero
supports are at distance at least $d$ and a cell has diameter
$2l<d$.

For each chosen $v_i$ use Lemma~\ref{si:lem-scale-cell-price} with
$Q_i=Q_0(64n)^p$ and radius $r_i=r/\|v_i\|_\sharp$. Zero pieces can
be dropped. At fixed cutoff choose one cofinal subsequence for
these finitely many cell normalizations, and write $E_i$ for the
limits. For every subset $J\subset\{1,\ldots,n\}$, chessboard
domination \cite[Theorem~C]{LP} gives
\[
 \left|\rho\!\left(\prod_{i\in J}
                 e^{z_i\mathcal O_a(v_i)-E_i(z_i)}\right)\right|
          \le \prod_{i\in J}Q_i,\qquad |z_i|<r_i.
\]
This follows by applying the bound to the unnormalized cell
functions at fixed $z$, then taking the limits of the scalar
normalizations. In finite volume use the finite chessboard
inequality and $E_{v_i,\ell}$ directly. The variables in this
product are independent complex coordinates. Cauchy's formula,
once per coordinate, bounds the product moment centered at
$E_i'(0)$ by $\prod_{i\in J}Q_i/r_i$. Applied in state $\nu$
with one coordinate it gives
$|\nu(\mathcal O_a(v_i))-E_i'(0)|\le Q_i/r_i$.
Expanding the product about these constants over all subsets
proves the bound $2^n\prod_i Q_i/r_i$ with the requested
centering.

Sum the last estimate over piece choices. The cost from the
pieces of each test is at most $C(R+1)^4\delta^{-5}$ times
its norm, so the bound is
\[
   [C(R+1)^4]^n n^{(p+5)n}\bar d^{-5n}\prod_i\|f_i\|_\sharp .
\]
It includes the coarse-cutoff case. Since $n!\ge(n/e)^n$, the
claim holds with $k=\lceil p+5\rceil$ and an adjusted constant.
Complex tests follow by real and imaginary parts and absorbing
$2^n$.

Cumulants of order $n\ge2$ do not change when constants are
subtracted. Use the moment formula on each block of a partition.
For block sizes $b_1,\ldots,b_m$ one has $\prod_j b_j!\le n!$.
After removing a factor $(n!)^{k-1}$, order the blocks and
sum over their positive sizes. At fixed ordered sizes there
are $n!/\prod_j b_j!$ ordered partitions. Hence the remaining
sum, including the coefficient $(m-1)!$, is
\[
    n!\sum_{m=1}^n \frac1m \binom{n-1}{m-1}\le 2^{n-1}n!.
\]
This proves the cumulant bound.
\end{proof}

For a version with a separation factor independent of $n$,
apply $x^n\le e^x n!$ with $x=\bar d^{-5}$ in
\eqref{si:eq-graded-moment}. This replaces $(n!)^k$ by
$(n!)^{k+1}$ and the separation factor by $e^{\bar d^{-5}}$.


\section{Uniform insertion in the reflected norm}
\label{si:sec-insertion}

Write $\tau_v$ for translation by $v$, so the support moves by
$v$, and $\theta_t$ for time reflection at $x_4=t$. At fixed
spacing $a$, choose a cell $C$ of side $l=L^j\ge a$ centered at
a lattice site, with lower time face $t$. The ratio $l/a$ is
odd. Its faces are link mirrors and its middle plane is a site
mirror. For a periodic Wilson Gibbs state $\rho$, the form
\begin{equation}\label{si:eq-reflection-form}
  \langle B,A\rangle_t=\rho(\overline{\theta_t B}A),
  \qquad \|A\|_t^2=\langle A,A\rangle_t
\end{equation}
is positive semidefinite by the all-mirror theorem
\cite[Theorem~B]{LP}. Its raw algebra consists of bounded
Borel functions of finitely many bonds whose endpoints lie
on the positive side of the link plane.

\begin{samepage}
\begin{lemma}[Cell multiplication and a time shift]
\label{si:lem-reflected-cell-operator}
For a bounded complex cell function $D$ using bonds with both
endpoints in $C$, and any cofinal sequence of compatible tori,
set $q_D=\limsup_{\ell}p_{l,\ell}(D)$. On the algebra of
\eqref{si:eq-reflection-form}, the map $T_D A=D\tau_{le_4}A$
satisfies
\begin{equation}\label{si:eq-cell-operator}
 \|T_D A\|_t\le q_D\|A\|_t,\qquad
 |\langle B,T_D A\rangle_t|\le q_D\|B\|_t\|A\|_t .
\end{equation}
It therefore descends to the quotient by the null space.
\end{lemma}
\end{samepage}

\begin{proof}
Put $D^*=\overline{\theta_{t+l/2}D}$. The reflection identities
$\theta_t\tau_{le_4}=\tau_{-le_4}\theta_t$ and
$\overline{\theta_tD^*}=\tau_{-le_4}D$, followed by translation
invariance of $\rho$, give
$\langle B,T_D A\rangle_t=\langle T_{D^*}B,A\rangle_t$.
These are identities on the raw algebra. For $K_D=T_{D^*}T_D$ and
$m_b=\langle A,K_D^b A\rangle_t$, they imply
\[
 m_{2s}=\|K_D^s A\|_t^2,\qquad
 m_{2s+1}=\|T_D K_D^s A\|_t^2 .
\]
Cauchy--Schwarz gives $m_b^2\le m_{b-1}m_{b+1}$ for $b\ge1$.
For $b=2s+1$ use $K_D^s A,K_D^{s+1} A$. For $b=2s\ge2$ use
$T_DK_D^{s-1}A,T_DK_D^sA$. If $m_0=0$, Cauchy--Schwarz gives
$m_1=0$. Otherwise $m_1^N\le m_0^{N-1}m_N$. In the nonzero
case this follows by iterating the increasing ratios
$m_b/m_{b-1}$.

In $m_N$ there are the two endpoint functions
$\overline{\theta_tA}$ and $\tau_{2Nle_4}A$, with $2N$ alternating
copies of $D^*,D$ between them. Choose a coarse side $l'=L^u l$
large enough to contain the finite spatial and time support of
$A$ in one coarse cell above $t$. Center the spatial directions
of that cell as in $C$ and set its lower time face at $t$. Its
faces are fine-board faces because $L^u$ is odd. Choose $N$ a
multiple of $L^u$. The two endpoints are then in separate
coarse cells outside the column. All used bonds of a constituent
fine cell lie in the open coarse cell containing it.

Apply fixed-spacing product domination on this coarse board
\cite[Theorem~C]{LP}, taking volume to infinity for the fixed
$N$. The two endpoint seminorms are at most $\|A\|_\infty$.
For each middle cell the fine-to-coarse estimate in that theorem
bounds its seminorm by the product of those of its fine factors.
Translations changing the reference cell preserve the fine
seminorm. Reflection through the middle of $C$ permutes the
fine board and conjugates its copy reflections, so invariance
of the torus expectation and conjugation invariance give
$p_{l,\ell}(D^*)=p_{l,\ell}(D)$. Since $q_D\le\|D\|_\infty$,
for every $\epsilon>0$ this proves
$m_N\le\|A\|_\infty^2(q_D+\epsilon)^{2N}$.
Take $N$th roots, send $N$ to infinity along the indicated
multiples, then let $\epsilon$ tend to zero. This includes
$q_D=0$ and gives the first estimate in
\eqref{si:eq-cell-operator}. The second is Cauchy--Schwarz.
\end{proof}

\begin{samepage}
\begin{theorem}[A centered lattice insertion]\label{si:thm-lattice-insertion}
Assume the ordinary source estimates and actual rooted domain
indexing of Lemma~\ref{si:lem-scale-cell-price}. Let $a\le l\le1$ be an
allowed cell side and $h$ a real Lipschitz test supported in the
radius-$l/4$ ball about its center. For any two periodic Wilson
Gibbs states $\rho,\nu$ at the same bare datum,
\begin{equation}\label{si:eq-lattice-insertion}
 \begin{aligned}
 &\left|\rho\!\left(\overline{\theta_tB}
    [\mathcal O_a(h)-\nu(\mathcal O_a(h))]\tau_{le_4}A\right)\right|\\
 &\hspace{25mm}\le C\|h\|_\sharp\|B\|_t\|A\|_t .
 \end{aligned}
\end{equation}
The constant is uniform in the states, cutoff, cell side and
the functions $A,B$ in the raw algebra.
\end{theorem}
\end{samepage}

\begin{proof}
For $a\le l/64$, Lemma~\ref{si:lem-scale-cell-price} with
$\delta=l/4$ yields one-variable normalizers with
\[
 p_{l,\ell}(e^{z\mathcal O_a(h)-E_{h,\ell}(z)})\le Q,
              \qquad |z|\|h\|_\sharp<r,
\]
with fixed $Q$. If $3\le l/a<64$,
every contributing plaquette is internal to $C$. Its vertices
in each coordinate are at distance at most $l/4+a/2<l/2$
from the cell center. The number of plaquettes is bounded by
$C64^4$. Take $E_{h,\ell}=0$ in this case and enlarge $Q$.
At $l/a=1$ the ball contains no plaquette centre, so the same
bound holds.

Take a cofinal subsequence at fixed cutoff with normalizer
limit $E_h$. For each $z$ in the disc, homogeneity and
convergence imply $q_{D_z}\le Q$ for
$D_z=e^{z\mathcal O_a(h)-E_h(z)}$. Apply
\eqref{si:eq-cell-operator} to $D_z$. Cauchy's formula applies
to the expectation because the fixed-cutoff observable is
bounded and $E_h$ is holomorphic. The derivative at zero is
bounded by $(Q/r)\|h\|_\sharp\|B\|_t\|A\|_t$ and inserts
$\mathcal O_a(h)-E_h'(0)$. With $A=B=1$ in state $\nu$ the
same argument bounds $|\nu(\mathcal O_a(h))-E_h'(0)|$ by
$(Q/r)\|h\|_\sharp$. The $D=1$ case of
\eqref{si:eq-cell-operator} bounds the resulting scalar term.
The same normalizer works in both states. If $\|h\|_\sharp=0$
the insertion vanishes directly.
\end{proof}

\begin{samepage}
\begin{proposition}[Insertion for already convergent separated moments]
\label{si:prop-limit-insertion}
Assume the hypotheses of Theorem~\ref{si:thm-lattice-insertion}
uniformly along $a_j\downarrow0$. Suppose that at each datum
$\rho_j,\nu_j$ are two periodic Wilson Gibbs states and the
centered moments converge on every fixed separated tuple of
compactly supported $C^1$ tests. Denote the limits by $\sigma$,
write $\phi(f)$ for an entry with test $f$, and set $\sigma(1)=1$.
Let $F,G$ be finite linear combinations of monomials whose
entries have mutually separated supports within each monomial
and are compactly supported $C^1$ tests lying strictly in positive
time. The form
$\langle G,F\rangle_\sigma=\sigma(\theta_0\overline G\,F)$
is positive semidefinite. For $l\le1$ an allowed power of $L$
and real $h\in C_c^1$ supported in the radius-$l/4$ ball
about $(y,0)$,
\begin{equation}\label{si:eq-limit-insertion}
 |\sigma(\theta_0\overline G\,
       \phi(\tau_{(l/2)e_4}h)\tau_{le_4}F)|
 \le C\|h\|_\sharp\|G\|_\sigma\|F\|_\sigma .
\end{equation}
Time translations by $u e_4$ for $u\ge0$ act by contractions
in this seminorm. Both operations respect null vectors.
\end{proposition}
\end{samepage}

\begin{proof}
Multilinear telescoping and Theorem~\ref{si:thm-graded-moments}
give equicontinuity on each fixed separated tuple when each
test is translated by a vector tending to zero. Indeed the
support of a difference is in the union of two nearby supports,
the gaps stay positive, and the difference tends to zero in
the sharp norm for compact $C^1$ tests. This also proves
translation invariance of the limits, by approximating a fixed
real shift with lattice shifts and using the invariance of
both states at each datum.

Let $y_j$ be a nearest spatial lattice point to $y$ and put
$v_j=(y_j-y,a_j/2)$ and $t_j=a_j/2$. The point
$(y_j,t_j+l/2)$ is a lattice site for all sufficiently fine
cutoffs. Translate $F,G$ by $v_j$ and the inserted test by
$v_j+(l/2)e_4$. The equality
$\theta_{t_j}\tau_{v_j}=\tau_{v_j}\theta_0$ identifies the
norm tuples. The centering state $\nu_j$ is invariant under
this lattice mirror by \cite[Theorem~B]{LP}, so its reflected
centering constants agree as well. For the lattice form apply
\eqref{si:eq-reflection-form} to the polynomials of centered
lattice observables. All supports are in the indicated open
half-spaces when $j$ is large, since the tuple gaps are fixed
and each plaquette has diameter $O(a_j)$. Positivity of the
limiting form follows. Apply \eqref{si:eq-lattice-insertion}
to the same polynomials and take the limit using the fixed-gap
equicontinuity. This gives \eqref{si:eq-limit-insertion} with
its stated constant even when a norm vanishes.

At fixed cutoff and $u\ge0$ a lattice time, invariance and
reflection Cauchy--Schwarz give
\[
 \|\tau_{ue_4}A\|_t^2\le\|A\|_t\|\tau_{2ue_4}A\|_t.
\]
If $\|A\|_t=0$ this gives the assertion directly. Otherwise
iterate, using $\|\tau_{2^kue_4}A\|_t\le\|A\|_\infty$, and
send $k$ to infinity. For any nonnegative real $u$, take
nonnegative lattice times tending to $u$ and use the same
fixed-gap passage. The norm bound for the insertion itself
follows by using the polynomial
$G=\phi(\tau_{(l/2)e_4}h)\tau_{le_4}F$
as the left argument of \eqref{si:eq-limit-insertion}.
This polynomial has separated
entries. The two norm estimates give null-space descent.
\end{proof}

\begin{samepage}
\begin{corollary}[Near-mirror insertion]\label{si:cor-mirror-insertion}
Under Proposition~\ref{si:prop-limit-insertion}, let
$0<\eta\le1$ and $f\in C_c^1$ be supported in
$\eta\le x_4\le3\eta$ with support diameter at most $R$. Then
\begin{equation}\label{si:eq-near-slab}
 \|\phi(f)\tau_{4\eta e_4}F\|_\sigma
 \le C(1+R)^3\eta^{-4}\|f\|_{C^1}\|F\|_\sigma .
\end{equation}
Fix $0<c<1/64$. If $k\in C_c^1$ is supported in the ball of
radius $c\eta$ about $y$ with $|y_4|\le\eta/8$, and the
entries of $F,G$ lie strictly above $\eta/4$, then
\begin{equation}\label{si:eq-near-mirror}
 |\sigma(\theta_0\overline F\,\phi(k)\,G)|
 \le C\eta^{-4}\|k\|_\sharp
     \|\tau_{-(\eta/4)e_4}F\|_\sigma
     \|\tau_{-(\eta/4)e_4}G\|_\sigma .
\end{equation}
The constants are independent of the number of entries in the
monomials and the number of monomials in $F,G$.
\end{corollary}
\end{samepage}

\begin{proof}
Take the largest power $l$ of $L$ at most $\eta$ and use a
smooth partition at mesh comparable to $l$ with pieces
supported in balls of radius $l/4$. There are at most
$C(1+R)^3\eta^{-3}$ nonzero pieces, since the support has
time width $2\eta$. The original test vanishes on the two
time faces, so $\|f\|_\infty\le\eta\|\partial_4 f\|_\infty$.
As $l>\eta/L$, the derivative of a partition factor costs
at most a fixed multiple of $\|f\|_{C^1}$ in the sharp norm
of the piece.

For a piece with center time $b$, let $t=b-l/2$.
Translate its product with $\theta_0\overline G$ and
$\tau_{4\eta e_4}F$ by $-t e_4$. The left polynomial
becomes the reflection of $\tau_{te_4}G$, and the right
one is $\tau_{le_4}\tau_{(4\eta-t-l)e_4}F$. Both extra
translations have times at least $\eta-3l/4>0$ because
$\eta-l/4\le b\le3\eta+l/4$. Apply
\eqref{si:eq-limit-insertion} and time contraction to each
piece. The resulting bilinear estimate has cost
$C(1+R)^3\eta^{-3}\|f\|_{C^1}\|G\|_\sigma\|F\|_\sigma$.
Set $G=\phi(f)\tau_{4\eta e_4}F$ and use $\eta^{-3}\le\eta^{-4}$.
Real and imaginary parts handle a complex $f$.

For \eqref{si:eq-near-mirror}, set $\xi=c\eta$ and translate
the whole expectation in time by $2\xi-y_4$. The translated
test lies in $[\xi,3\xi]$ in time. The left polynomial is
the reflection of $\tau_{(y_4-2\xi)e_4}F$, and the right
polynomial is $\tau_{4\xi e_4}\tau_{(-y_4-2\xi)e_4}G$.
The two polynomials before reflection and the $4\xi$ shift
are positive translates of their $\eta/4$ downshifts, by
at least $\eta/8-2\xi>0$. Apply \eqref{si:eq-near-slab},
Cauchy--Schwarz and time contraction. The translated test
has diameter at most $2c$, and the factor $c^{-4}$ is
fixed. This gives the asserted bound.
\end{proof}


\section{A two-source estimate at running stops}\label{si:sec-running-stop}

The preceding transfer can be used at a finer first-passage stop
than the unit-scale stop. The needed ordinary estimates must be
available at each of these stops, with their own site counts.

\begin{samepage}
\begin{proposition}[Running-stop estimate for separated pairs]
\label{si:prop-running-stop-mirror}
Take first-passage data of \cite{RT} with terminal index $K$ and
reference floor $h$. Assume the self-calibrated prefix is common
under the admitted changes of units. On the cofinal tori of every
first-passage stop used below, assume the ordinary source bounds
in the units of that stop, with $C_++C_-\le C(1+\log x)$ at
terminal weight $x$. The source radius, local energy bounds,
entropy margin and geometric constants are uniform over these
stops. The minimum allowed torus side may depend on the stop.

Fix $\vartheta>0$. Suppose $f_1,f_2$ have support radii at most $d$,
mutual support distance at least $d$, and
$\operatorname{Lip} f_i\le\|f_i\|_\infty/(\vartheta d)$.
There are $c,C>0$ such that, for $0<d\le c$ and every periodic
Wilson state $\rho$ at the original datum,
\begin{equation}\label{si:eq-running-stop-mirror}
 \big|\kappa_2^\rho(\mathcal O_a(f_1),\mathcal O_a(f_2))\big|
 \le C\big[1+h+\log(1/d)\big]^C
                            \|f_1\|_\infty\|f_2\|_\infty .
\end{equation}
The constants are uniform in cutoff and support positions. The
bound holds in fixed-gap limits of these states whenever the
separated second cumulant converges. In particular multiplication
of the right side by $d^2$ tends to zero as $d\downarrow0$ at
fixed $h$.
\end{proposition}
\end{samepage}
\begin{proof}
In the real trajectory, call $i$ a record if $x_i<x_t$ for all
$t<i$. The initial and first-passage indices are records. Let
$J=1+\lceil c_5/\lambda_{\rm s}\rceil$, with the lower drift
constants of \cite[Theorem~A and Proposition~4.4]{RT}. After a record $i$ with $i+J\le K$,
$x_i-x_{i+J}\ge\lambda_{\rm s}J-c_5>0$, so there is a new record
within $J$ steps. If fewer steps remain, $K$ supplies the next one.
Thus the depths of the records, listed from $0$ to $K$, have gaps
at most $J$.

Put $n=\lfloor\log_L(c/d)\rfloor$. For $K<n$, the original spacing
is greater than $d/c$. A support of radius $d$ then contains at
most $C(1+d/a)^4$ plaquette barycentres, and the elementary bound
$0\le a_p\le2$ proves \eqref{si:eq-running-stop-mirror}. If $K\ge n$,
choose the largest record depth $b\le n$. Then
$c L^{-J-1}\le L^b d\le c$. For $b>0$, the prefix ending at $K-b$
is a first-passage trajectory for the new floor
$h_b=x_{K-b}-B_{\rm s}$, where $B_{\rm s}$ is the original upper
step constant. Every earlier entry exceeds $x_{K-b}$ by the
record property. The old first-passage condition gives $h_b>h$.
The upper drift and the terminal window give
$x_{K-b}\le h+B_{\rm s}(b+1)$. For $b=0$ keep the original stop.

Dilate by $L^b$. The bijection of the bare plaquettes preserves
both Wilson conditional specifications and their scalar
observables with tests $f_i^{[b]}(x)=f_i(L^{-b}x)$. Hence it carries
periodic Gibbs states to periodic Gibbs states and preserves the
cumulant. The common-prefix hypothesis identifies the stopped
normalization in the new units. Now the support radii are at most
$c$, the distance is bounded below by $c L^{-J-1}$, and
$\|f_i^{[b]}\|_\sharp\le C\|f_i\|_\infty$.

Apply \cref{si:thm-state-transfer} in these units at the new stop.
Its proof uses at most two occupied cells of bounded side. The
cell normalizer bound of \cref{si:lem-cell-normalizer} is the sum
of $(C_++C_-)$ times that bounded cell volume and a uniform
energy tail. Therefore the residual generating function has
bound $Q\le\exp[C(1+\log x_{K-b})]$ on a fixed source disc.
The mixed derivative of the common holomorphic energy is bounded
uniformly at the new separation floor by the local-energy
estimate in \cref{si:cor-state-cumulants}. For the residual $H$
with $H(0)=1$, its log derivative is
$\partial_{12}H(0)-\partial_1H(0)\partial_2H(0)$. Cauchy gives
a bound $C(Q+Q^2)\prod_i\|f_i^{[b]}\|_\sharp$. Inserting the
upper bound on $x_{K-b}$ proves \eqref{si:eq-running-stop-mirror}.
The volume limit in the transfer is at this fixed spacing and
stop, so a stop-dependent minimum torus side causes no loss.
The fixed-gap passage follows by convergence of the cumulant.
\end{proof}

\begin{thebibliography}{Bal88c}
\small
\bibitem{BA} \emph{Block averaging and gauge geometry for four-dimensional lattice Yang--Mills theory}, companion manuscript, 2026.
\bibitem{LP} \emph{Periodic Wilson states and reflection positivity}, companion manuscript, 2026.
\bibitem{EM} \emph{Iteration of regular interactions on prescribed coupling prefixes}, companion manuscript, 2026.
\bibitem{RT} \emph{Coupling trajectories and partition stability in lattice Yang--Mills theory}, companion manuscript, 2026.
\bibitem{Bal88c} T.~Ba{\l}aban, \emph{Renormalization group approach to lattice gauge field theories. II. Cluster expansions}, Comm. Math. Phys. \textbf{116} (1988), 1--22. \href{https://doi.org/10.1007/BF01239022}{doi:10.1007/BF01239022}.
\end{thebibliography}
\end{document}
